\chapter{Dirac 与电磁场的量子化}\label{chap:qed}

量子场论的通用结构现在落实到电子与电磁场：电子由 Dirac 场描述，电磁场由 Maxwell 四势描述，局域 $U(1)$ 对称性通过最小耦合把二者连接。Dirac 场的正、负频分支分别组织电子和正电子自由度；Maxwell 四势的四个分量则受到规范约束，只有两个横向偏振成为独立的物理光子。

量子光学已经在 Coulomb 规范的横向模式上算出原子偶极辐射。这里进一步保留相对论费米子与规范约束，使自由传播子、协变规范和外态归一化可以用于电子、正电子与光子的散射。\cref{chap:electromagnetic-scattering}将从同一电流顶角计算过程和辐射修正；\cref{chap:quantum-electromagnetic-media}再把材料、边界和损耗加入电磁响应。

\section{Dirac 与 Maxwell 场量子化}

一般量子场论已经说明怎样从自由模得到 Fock 空间、传播子和散射振幅。电磁相互作用只需指定两种具体场：带电 Dirac 场承载电子与正电子，Maxwell 四势承载光子；两者由最小耦合连接。Maxwell 四势还含规范冗余，因此量子化时必须同时保留约束结构，最终只让两个横向偏振进入物理光子空间。

\subsection{Dirac 场量子化}

\begin{definition}[自由 Dirac 场的模展开与费米代数]
自由 Dirac 方程具有正频和负频两支质量壳解。量子理论需要同时满足正定内积、能谱下界和局域等时反对易关系；这三项条件把正频系数解释为电子湮灭算符 $b_s$，把负频系数解释为正能正电子的产生算符 $d_s^\dagger$。采用普通三维 Dirac--$\delta$ 的振子归一化时，场展开系数还必须补偿自旋子的协变归一化，因此
\begin{equation}
\boxed{
\hat\psi(x)
=\sum_s\int\frac{\dd^3 p}{(2\pi\hbar)^3}\sqrt{\frac{c}{2E_{\mathbf p}}}
\left[
\hat b_s(\mathbf p)u_s(\mathbf p)\ee^{-ip\cdot x/\hbar}
+\hat d_s^\dagger(\mathbf p)v_s(\mathbf p)\ee^{+ip\cdot x/\hbar}
\right].}
\label{eq:dirac-field-exp-r2}
\end{equation}
其中 $p^0=E_{\mathbf p}/c>0$，并采用 $u_s^\dagger u_{s'}=v_s^\dagger v_{s'}=(2E_{\mathbf p}/c)\delta_{ss'}$ 的自旋子归一化。系数 $\sqrt{c/(2E_{\mathbf p})}$ 与普通 $\delta^{(3)}$ 振子代数配合，保证等时场反对易关系具有单位系数。

费米统计由正则反对易代数实现：
\begin{equation}
\boxed{
\{b_r(\mathbf p),b_s^\dagger(\mathbf q)\}
=\{d_r(\mathbf p),d_s^\dagger(\mathbf q)\}
=(2\pi\hbar)^3\delta_{rs}\delta^{(3)}(\mathbf p-\mathbf q),}
\label{eq:fermion-anticomm-normalization-r2}
\end{equation}
其余反对易子为零。连续场与振子都是算符分布，算符等式须以测试函数涂抹后理解。
\end{definition}

\begin{proposition}[Dirac 场的模式 CAR 归一化]
假设正文的连续旋量完备关系与等时场 CAR；系数由相同动量测度固定。则有以下结论（原式 \eqnrefs{eq:fermion-anticomm-normalization-r2}）：
\begin{equation*}
\boxed{
\{b_r(\mathbf p),b_s^\dagger(\mathbf q)\}
=\{d_r(\mathbf p),d_s^\dagger(\mathbf q)\}
=(2\pi\hbar)^3\delta_{rs}\delta^{(3)}(\mathbf p-\mathbf q),}
\end{equation*}
\end{proposition}
\begin{proof}
先由等时场 CAR 提取模式算符。记固定时刻的空间 Fourier 变换为
\[
\widetilde\psi(t,\mathbf p)
=\int\dd^3x\,\ee^{-\ii\mathbf p\cdot\mathbf x/\hbar}\psi(t,\mathbf x).
\]
自由模式展开及 $E_{-\mathbf p}=E_{\mathbf p}$ 给出
\[
\widetilde\psi(t,\mathbf p)
=\sqrt{\frac{c}{2E_{\mathbf p}}}
\sum_s\left[
b_s(\mathbf p)u_s(\mathbf p)\ee^{-\ii E_{\mathbf p}t/\hbar}
+d_s^\dagger(-\mathbf p)v_s(-\mathbf p)\ee^{+\ii E_{\mathbf p}t/\hbar}
\right].
\]
同一 Hamiltonian 动量的正负能基满足
\[
u_r^\dagger(\mathbf p)u_s(\mathbf p)
=v_r^\dagger(\mathbf p)v_s(\mathbf p)
=\frac{2E_{\mathbf p}}c\delta_{rs},
\qquad
u_r^\dagger(\mathbf p)v_s(-\mathbf p)=0.
\]
因此模式展开有精确逆变换
\begin{align*}
b_s(\mathbf p)
&=\sqrt{\frac{c}{2E_{\mathbf p}}}\,
\ee^{+\ii E_{\mathbf p}t/\hbar}
u_s^\dagger(\mathbf p)\widetilde\psi(t,\mathbf p),\\
d_s^\dagger(\mathbf p)
&=\sqrt{\frac{c}{2E_{\mathbf p}}}\,
\ee^{-\ii E_{\mathbf p}t/\hbar}
v_s^\dagger(\mathbf p)\widetilde\psi(t,-\mathbf p).
\end{align*}
等时场 CAR 的 Fourier 变换为
\begin{align*}
\{\widetilde\psi_\alpha(t,\mathbf p),
\widetilde\psi_\beta^\dagger(t,\mathbf q)\}
&=(2\pi\hbar)^3\delta^{(3)}(\mathbf p-\mathbf q)\delta_{\alpha\beta},\\
\{\widetilde\psi_\alpha(t,\mathbf p),
\widetilde\psi_\beta(t,\mathbf q)\}&=0.
\end{align*}
把逆变换代入，先保留两个独立动量，得到
\begin{align*}
\{b_r(\mathbf p),b_s^\dagger(\mathbf q)\}
&=\frac{c\,\ee^{\ii(E_{\mathbf p}-E_{\mathbf q})t/\hbar}}
{2\sqrt{E_{\mathbf p}E_{\mathbf q}}}
u_r^\dagger(\mathbf p)u_s(\mathbf q)
(2\pi\hbar)^3\delta^{(3)}(\mathbf p-\mathbf q)\\
&=(2\pi\hbar)^3\delta_{rs}\delta^{(3)}(\mathbf p-\mathbf q),\\
\{d_r^\dagger(\mathbf p),d_s(\mathbf q)\}
&=\frac{c\,\ee^{-\ii(E_{\mathbf p}-E_{\mathbf q})t/\hbar}}
{2\sqrt{E_{\mathbf p}E_{\mathbf q}}}
v_r^\dagger(\mathbf p)v_s(\mathbf q)
(2\pi\hbar)^3\delta^{(3)}(\mathbf p-\mathbf q)\\
&=(2\pi\hbar)^3\delta_{rs}\delta^{(3)}(\mathbf p-\mathbf q).
\end{align*}
第二式利用反对易子的对称性及交换哑指标，即为所述 $\{d_r,d_s^\dagger\}$。
$\{b,d^\dagger\}$、$\{b,b\}$、$\{d^\dagger,d^\dagger\}$ 由第二条 Fourier CAR 消失；
$\{b_r(\mathbf p),d_s(\mathbf q)\}$ 则含
$\delta^{(3)}(\mathbf p+\mathbf q)u_r^\dagger(\mathbf p)v_s(\mathbf q)$，
由反向动量的正交关系为零。取伴随覆盖其余组合。
这直接从场 CAR 证明了\eqnrefs{eq:fermion-anticomm-normalization-r2}。

下面完整核对反向关系：由模式 CAR 恢复等时场 CAR，并检查模式测度与旋量归一化的一致性。
把正文的自由场展开写成

\begin{align*}
\psi_\alpha(t,\mathbf x)
&=\sum_s\int\frac{\dd^3 p}{(2\pi\hbar)^3}
\sqrt{\frac{c}{2E_{\mathbf p}}}
\left[
 b_s(\mathbf p)u_{s,\alpha}(\mathbf p)\ee^{-\ii p\cdot x/\hbar}
 +d_s^\dagger(\mathbf p)v_{s,\alpha}(\mathbf p)\ee^{+\ii p\cdot x/\hbar}
\right],\\
\psi_\beta^\dagger(t,\mathbf y)
&=\sum_{s'}\int\frac{\dd^3 p'}{(2\pi\hbar)^3}
\sqrt{\frac{c}{2E_{\mathbf p'}}}
\left[
 b_{s'}^\dagger(\mathbf p')u_{s',\beta}^\dagger(\mathbf p')\ee^{+\ii p'\cdot y/\hbar}
 +d_{s'}(\mathbf p')v_{s',\beta}^\dagger(\mathbf p')\ee^{-\ii p'\cdot y/\hbar}
\right].
\end{align*}
只有 $\{b,b^\dagger\}$ 与 $\{d,d^\dagger\}$ 两类项保留。用模算符反对易关系先完成 $\mathbf p'$ 积分，在等时 $x^0=y^0$ 得
\begin{align*}
\{\psi_\alpha(t,\mathbf x),\psi_\beta^\dagger(t,\mathbf y)\}
&=\int\frac{\dd^3 p}{(2\pi\hbar)^3}\frac{c}{2E_{\mathbf p}}
\sum_s\Bigl[
 u_{s,\alpha}(\mathbf p)u_{s,\beta}^\dagger(\mathbf p)\ee^{+\ii\mathbf p\cdot(\mathbf x-\mathbf y)/\hbar}\\
&\hspace{7em}+v_{s,\alpha}(\mathbf p)v_{s,\beta}^\dagger(\mathbf p)\ee^{-\ii\mathbf p\cdot(\mathbf x-\mathbf y)/\hbar}
\Bigr].
\end{align*}
第二项作变量替换 $\mathbf p\mapsto-\mathbf p$，于是两项具有同一 Fourier 相位：
\begin{align*}
\{\psi_\alpha(t,\mathbf x),\psi_\beta^\dagger(t,\mathbf y)\}
&=\int\frac{\dd^3 p}{(2\pi\hbar)^3}\frac{c}{2E_{\mathbf p}}
\left[
\sum_su_s(\mathbf p)u_s^\dagger(\mathbf p)
+\sum_sv_s(-\mathbf p)v_s^\dagger(-\mathbf p)
\right]_{\alpha\beta}
 \ee^{+\ii\mathbf p\cdot(\mathbf x-\mathbf y)/\hbar}.
\end{align*}
正、负能旋量完备关系分别为
\begin{equation*}
\sum_su_s(\mathbf p)u_s^\dagger(\mathbf p)=(\slashed p+m_{\mathrm e} c)\gamma^0,
\qquad
\sum_sv_s(-\mathbf p)v_s^\dagger(-\mathbf p)=(\slashed p'-m_{\mathrm e} c)\gamma^0,
\end{equation*}
其中 $p=(E_{\mathbf p}/c,\mathbf p)$、$p'=(E_{\mathbf p}/c,-\mathbf p)$。两式相加时空间动量项与质量项分别抵消，只剩
\begin{align*}
(\slashed p+m_{\mathrm e} c)\gamma^0+(\slashed p'-m_{\mathrm e} c)\gamma^0
&=(\slashed p+\slashed p')\gamma^0\\
&=\frac{2E_{\mathbf p}}{c}\,\gamma^0\gamma^0
=\frac{2E_{\mathbf p}}{c}\,\id(4).
\end{align*}
因此展开中的 $c/(2E_{\mathbf p})$ 恰好相消，最后用 Fourier 表示
\begin{equation*}
\int\frac{\dd^3 p}{(2\pi\hbar)^3}\ee^{+\ii\mathbf p\cdot(\mathbf x-\mathbf y)/\hbar}
=\delta^{(3)}(\mathbf x-\mathbf y)
\end{equation*}
得到
\begin{equation*}
\{\psi_\alpha(t,\mathbf x),\psi_\beta^\dagger(t,\mathbf y)\}
=\delta_{\alpha\beta}\delta^{(3)}(\mathbf x-\mathbf y),
\end{equation*}
这恢复了等时场 CAR；与前面的逆变换计算合并，得到场与模式两组正则关系的等价性。
\end{proof}


\begin{proposition}[正则归一化与正能粒子谱]
上述模展开和反对易关系给出等时正则关系
\begin{equation*}
\{\psi_\alpha(\mathbf x,t),\psi_\beta^\dagger(\mathbf y,t)\}
=\delta_{\alpha\beta}\delta^{(3)}(\mathbf x-\mathbf y).
\end{equation*}

正规序后的自由哈密顿量因而为
\begin{equation}
\boxed{
:\!H_D\!:
=\sum_s\int\frac{\dd^3 p}{(2\pi\hbar)^3}E_{\mathbf p}
\left[b_s^\dagger(\mathbf p)b_s(\mathbf p)
+d_s^\dagger(\mathbf p)d_s(\mathbf p)\right].}
\label{eq:dirac-H-positive-r2}
\end{equation}
在有限盒的单模记号中，同一反对易代数给 $(b_a^\dagger)^2=(d_a^\dagger)^2=0$；连续表象对应任意同一归一化波包产生算符的平方为零。因此费米 Fock 空间由反对称张量幂构成，Pauli 不相容性与正能谱由同一代数结构共同决定。
\end{proposition}

\begin{proposition}[自由 Dirac 二次算子的正规序谱]
采用自由 Dirac 模式完备关系和 CAR，并将真空常数按同一正规序处方消去；结论为自由模式上的算子恒等式。则有以下结论（原式 \eqnrefs{eq:dirac-H-positive-r2}）：
\begin{equation*}
\boxed{
:\!H_D\!:
=\sum_s\int\frac{\dd^3 p}{(2\pi\hbar)^3}E_{\mathbf p}
\left[b_s^\dagger(\mathbf p)b_s(\mathbf p)
+d_s^\dagger(\mathbf p)d_s(\mathbf p)\right].}
\end{equation*}
\end{proposition}
\begin{proof}
自由 Dirac 单粒子 Hamiltonian 为
\[
 \hat h_D=-\ii\hbar c\,\boldsymbol\alpha\cdot\boldsymbol\nabla
 +\beta_Dm_{\mathrm e} c^2.
\]
先看场展开中两类空间--时间模本身。正频模
\[
 \Phi^{(+)}_{\mathbf p s}(x)
 =u_s(\mathbf p)\ee^{-\ii p\cdot x/\hbar}
\]
满足 $\ii\hbar\partial_t\Phi^{(+)}=+E_{\mathbf p}\Phi^{(+)}$；由自由 Dirac 方程
$\ii\hbar\partial_t\Phi^{(+)}=\hat h_D\Phi^{(+)}$，所以它是 $\hat h_D$ 的 $+E_{\mathbf p}$ 本征模。负频模
\[
 \Phi^{(-)}_{\mathbf p s}(x)
 =v_s(\mathbf p)\ee^{+\ii p\cdot x/\hbar}
\]
则满足
\[
 \ii\hbar\partial_t\Phi^{(-)}_{\mathbf p s}
 =-E_{\mathbf p}\Phi^{(-)}_{\mathbf p s}
 =\hat h_D\Phi^{(-)}_{\mathbf p s}.
\]
因此正文场展开中的 $b_s$ 乘在 $+E$ 单粒子本征模上，而 $d_s^\dagger$ 乘在 $-E$ 单粒子本征模上；后一个负号正是未正规序 Hamiltonian 中
$-d_sd_s^\dagger$ 的来源。

为避免连续归一化中的 $\delta^{(3)}(\mathbf0)$ 干扰计数，先把体系放在有限盒中。用正交归一的单粒子模记号写
\[
 \hat\psi
 =\sum_a\left(
 \hat b_a\Phi_a^{(+)}+\hat d_a^\dagger\Phi_a^{(-)}
 \right),
\]
其中 $a=(\mathbf p,s)$。自由 Hamiltonian 是一体算符的二次量子化：
\begin{align*}
 \hat H_D
 &=\int\dd^3x\,\hat\psi^\dagger\hat h_D\hat\psi\\
 &=\sum_{a,b}\Bigl[
 b_a^\dagger b_b\langle\Phi_a^{(+)}|\hat h_D|\Phi_b^{(+)}\rangle
 +b_a^\dagger d_b^\dagger\langle\Phi_a^{(+)}|\hat h_D|\Phi_b^{(-)}\rangle\\
 &\hspace{30mm}
 +d_a b_b\langle\Phi_a^{(-)}|\hat h_D|\Phi_b^{(+)}\rangle
 +d_a d_b^\dagger\langle\Phi_a^{(-)}|\hat h_D|\Phi_b^{(-)}\rangle
 \Bigr].
\end{align*}
由于 $\hat h_D$ 厄米，不同本征值的模正交；同一动量处 $+E_{\mathbf p}$ 与 $-E_{\mathbf p}$ 不同，所以两个交叉矩阵元都为零。其余两个矩阵元为
\begin{align*}
 \langle\Phi_a^{(+)}|\hat h_D|\Phi_b^{(+)}\rangle
 &=+E_a\delta_{ab},\\
 \langle\Phi_a^{(-)}|\hat h_D|\Phi_b^{(-)}\rangle
 &=-E_a\delta_{ab}.
\end{align*}
于是有限盒中严格得到
\begin{equation*}
 \hat H_D=\sum_a E_a
 \left(b_a^\dagger b_a-d_a d_a^\dagger\right).
\end{equation*}
正文的因子 $\sqrt{c/(2E_{\mathbf p})}$ 与
$u_s^\dagger u_{s'}=v_s^\dagger v_{s'}=(2E_{\mathbf p}/c)\delta_{ss'}$
恰好把上述模内积化成 Kronecker--$\delta$；取无限体积极限后
$\sum_{\mathbf p}/V$ 恢复为 $\int\dd^3p/(2\pi\hbar)^3$。

现在才使用费米反对易关系。有限盒中
$\{d_a,d_b^\dagger\}=\delta_{ab}$，故逐模有
\[
 -d_ad_a^\dagger=d_a^\dagger d_a-1.
\]
因此
\begin{align*}
 \hat H_D
 &=\sum_aE_a\left(b_a^\dagger b_a+d_a^\dagger d_a\right)
 -\sum_aE_a.
\end{align*}
最后一项与占据态无关，是自由场真空零点常数；正规序把它减去。恢复连续动量归一化后便得到正文
\eqnrefs{eq:dirac-H-positive-r2}。所以“正电子具有正能量”并不是把经典负能解手工改号：负频模先给出 $-d d^\dagger$，再由费米反对易代数把它重排为 $+d^\dagger d$ 加真空常数。
\end{proof}


正文结果的两个归一化核验都在章末展开：等时反对易子中负频项必须先作 $\mathbf p\mapsto-\mathbf p$，再与正频项合并；Hamiltonian 则先保留 $-d d^\dagger$，再用费米重排得到 $+d^\dagger d$ 与真空常数。\eqnrefs{eq:feqo-hamiltonian-energy-projectors} 的 $\Pi_++\Pi_-=I_4$ 正是第一步计算的有限维内容。

\begin{exbox}[Dirac 场的错误统计]
量子化规则可以由理论必须同时具有能谱下界与正定态范数来检验。把模展开代入尚未正规序的自由哈密顿量时，电子和正电子部分具有结构
\begin{equation*}
 H_D\supset
 \sum_s\int\frac{\dd^3 p}{(2\pi\hbar)^3}E_{\mathbf p}
 \bigl(b_s^\dagger b_s-d_s d_s^\dagger\bigr).
\end{equation*}
如果错误地令正电子算符满足玻色关系 $[d,d^\dagger]=+1$，重排第二项得到
\begin{equation*}
 -d d^\dagger=-d^\dagger d-1,
\end{equation*}
于是每增加一个正电子，正规序后的能量反而减少 $E_{\mathbf p}$；不断增加占据数便使能谱没有下界。若改取 $[d,d^\dagger]=-1$ 试图修复能量符号，则一正电子态 $d^\dagger|0\rangle$ 的范数变为负值。两条路都不可接受。

反对易关系则给
\begin{equation*}
 \{d,d^\dagger\}=1
 \quad\Longrightarrow\quad
 -d d^\dagger=d^\dagger d-1,
\end{equation*}
所以电子和正电子都以 $+E_{\mathbf p}$ 计入激发能，同时 $(d^\dagger)^2=0$。正能谱、正定 Hilbert 空间和 Pauli 不相容性因此是同一反对易代数的三个后果；再结合 Noether 电荷便能同时确定 $d^\dagger$ 所产生激发的电荷符号。
\end{exbox}

全局 $U(1)$ Noether 电荷在正规序后成为
\begin{equation}
\boxed{
Q=q_{\mathrm e}\sum_s\int\frac{\dd^3 p}{(2\pi\hbar)^3}
\left[b_s^\dagger b_s-d_s^\dagger d_s\right].}
\label{eq:dirac-charge-r2}
\end{equation}

\begin{proposition}[自由 Dirac 电荷在 Fock 基上的作用]
采用正文电荷定义、模式归一化和正规序；共用稠密有限粒子检验域。则有以下结论（原式 \eqnrefs{eq:dirac-charge-r2}）：
\begin{equation*}
\boxed{
Q=q_{\mathrm e}\sum_s\int\frac{\dd^3 p}{(2\pi\hbar)^3}
\left[b_s^\dagger b_s-d_s^\dagger d_s\right].}
\end{equation*}
\end{proposition}
\begin{proof}
从

\begin{equation*}
Q=q_{\mathrm e}\int\dd^3 x\,\psi^\dagger(t,\mathbf x)\psi(t,\mathbf x)
\end{equation*}
出发，把正、负频模展开逐项相乘。空间积分给
\begin{equation*}
\int\dd^3 x\,\ee^{\ii(\mathbf p-\mathbf p')\cdot\mathbf x/\hbar}
=(2\pi\hbar)^3\delta^{(3)}(\mathbf p-\mathbf p'),
\end{equation*}
因此同频项先化为
\begin{align*}
Q_{bb}
&=q_{\mathrm e}\sum_{ss'}\int\frac{\dd^3 p}{(2\pi\hbar)^3}
\frac{c}{2E_{\mathbf p}}
 b_s^\dagger(\mathbf p)b_{s'}(\mathbf p)
 u_s^\dagger(\mathbf p)u_{s'}(\mathbf p),\\
Q_{dd}
&=q_{\mathrm e}\sum_{ss'}\int\frac{\dd^3 p}{(2\pi\hbar)^3}
\frac{c}{2E_{\mathbf p}}
 d_s(\mathbf p)d_{s'}^\dagger(\mathbf p)
 v_s^\dagger(\mathbf p)v_{s'}(\mathbf p).
\end{align*}
用旋量归一化
\begin{equation*}
u_s^\dagger u_{s'}=v_s^\dagger v_{s'}=\frac{2E_{\mathbf p}}{c}\delta_{ss'}
\end{equation*}
后，$c/(2E_{\mathbf p})$ 与旋量范数相消。正负频交叉项的空间积分给 $\delta^{(3)}(\mathbf p+\mathbf p')$，因此实际内积是 $u^\dagger(\mathbf p)v(-\mathbf p)$ 及其共轭。它们分别属于同一 $H_0(\mathbf p)$ 的 $+E$ 和 $-E$ 本征空间，由\eqnrefs{eq:feqo-negative-hamiltonian-spinor} 的正交性为零；不能把 $u^\dagger(\mathbf p)v(\mathbf p)$ 无条件设为零。于是
\begin{equation*}
Q=q_{\mathrm e}\sum_s\int\frac{\dd^3 p}{(2\pi\hbar)^3}
\left[b_s^\dagger(\mathbf p)b_s(\mathbf p)+d_s(\mathbf p)d_s^\dagger(\mathbf p)\right].
\end{equation*}
再使用反对易关系
\begin{align*}
d_s(\mathbf p)d_s^\dagger(\mathbf p)
&=\{d_s(\mathbf p),d_s^\dagger(\mathbf p)\}-d_s^\dagger(\mathbf p)d_s(\mathbf p)\\
&=\text{真空常数}-d_s^\dagger(\mathbf p)d_s(\mathbf p).
\end{align*}
按正文采用的正规序约定减去真空常数项后，得到
\begin{equation*}
:Q:=q_{\mathrm e}\sum_s\int\frac{\dd^3 p}{(2\pi\hbar)^3}
\left[b_s^\dagger(\mathbf p)b_s(\mathbf p)-d_s^\dagger(\mathbf p)d_s(\mathbf p)\right],
\end{equation*}
即\eqnrefs{eq:dirac-charge-r2}。作用在一粒子态上可直接检验
\begin{equation*}
:Q:\,b_s^\dagger|0\rangle=q_{\mathrm e}\,b_s^\dagger|0\rangle,
\qquad
:Q:\,d_s^\dagger|0\rangle=-q_{\mathrm e}\,d_s^\dagger|0\rangle.
\end{equation*}
由于 $q_{\mathrm e}<0$，$d_s^\dagger$ 产生的反粒子携带正电荷 $-q_{\mathrm e}$。
\end{proof}

因此 $d^\dagger$ 激发携带与电子相反的电荷。负频经典解在量子化后对应正能反粒子产生算符；Dirac 海是这一场论结构的历史单粒子表述。

\subsection{Dirac 传播子与因果性}

\begin{proposition}[Dirac 反对易核的因果支撑]
自由 Dirac 场的局域性由同一质量壳上的正、负频分支共同保证。为把这一事实与标量场的因果结构直接联系起来，定义标量 Pauli--Jordan 核
\begin{equation*}
 \Delta(x-y)
 \equiv
 \int\frac{\dd^3 p}{(2\pi\hbar)^3}\frac{c}{2E_{\mathbf p}}
 \left[
 \ee^{-\ii p\cdot(x-y)/\hbar}
 -\ee^{+\ii p\cdot(x-y)/\hbar}
 \right].
\end{equation*}
利用 Dirac 自旋子的完备关系，场反对易子可写成
\begin{equation}
 \boxed{
 \{\hat\psi_\alpha(x),\bar{\hat\psi}_\beta(y)\}
 =\left(\ii\hbar\slashed\partial_x+m_{\mathrm e} c\right)_{\alpha\beta}
 \Delta(x-y).}
 \label{eq:dirac-PJ-factor-r8}
\end{equation}
标量场的 Pauli--Jordan 核满足：当 $(x-y)^2<0$ 时，$\Delta(x-y)=0$。有限阶局域微分算符不会扩大分布的支撑，因此\eqnrefs{eq:dirac-PJ-factor-r8} 在类空间隔同样为零。费米场本身不是玻色型可观测量；由偶数个费米场构成的局域可观测量，例如 $\bar\psi\psi$ 与 $j^\mu=q_{\mathrm e} c\bar\psi\gamma^\mu\psi$，由这一反对易结构得到类空间隔处的对易性，从而满足微观因果性。

\end{proposition}

\begin{proposition}[Dirac 对易核的标量因式分解]
设自由 Dirac 算子与标量质量壳核采用相同 Fourier 和 CAR 约定；恒等式按算符值分布理解。则有以下结论（原式 \eqnrefs{eq:dirac-PJ-factor-r8}）：
\begin{equation*}
 \boxed{
 \{\hat\psi_\alpha(x),\bar{\hat\psi}_\beta(y)\}
 =\left(\ii\hbar\slashed\partial_x+m_{\mathrm e} c\right)_{\alpha\beta}
 \Delta(x-y).}
\end{equation*}
\end{proposition}
\begin{proof}
由\eqnrefs{eq:dirac-field-exp-r2}，$\hat\psi$ 含 $bu$ 与 $d^\dagger v$ 两类项，而 $\bar{\hat\psi}$ 含 $b^\dagger\bar u$ 与 $d\bar v$。四种交叉组合中只有

\begin{equation*}
 \{b_s(\mathbf p),b_{s'}^\dagger(\mathbf q)\},
 \qquad
 \{d_s^\dagger(\mathbf p),d_{s'}(\mathbf q)\}
\end{equation*}
非零。完成 $\mathbf q$ 积分并使用自旋完备关系后，得到
\begin{align*}
 \{\hat\psi_\alpha(x),\bar{\hat\psi}_\beta(y)\}
 &=\int\frac{\dd^3 p}{(2\pi\hbar)^3}\frac{c}{2E_{\mathbf p}}
 \bigg[
 (\slashed p+m_{\mathrm e} c)_{\alpha\beta}\ee^{-\ii p\cdot(x-y)/\hbar}\\
 &\hspace{42mm}
 +(\slashed p-m_{\mathrm e} c)_{\alpha\beta}\ee^{+\ii p\cdot(x-y)/\hbar}
 \bigg].
\end{align*}
另一方面，逐项作用 $\ii\hbar\slashed\partial_x+m_{\mathrm e} c$ 有
\begin{align*}
 (\ii\hbar\slashed\partial_x+m_{\mathrm e} c)\ee^{-\ii p\cdot(x-y)/\hbar}
 &=(\slashed p+m_{\mathrm e} c)\ee^{-\ii p\cdot(x-y)/\hbar},\\
 (\ii\hbar\slashed\partial_x+m_{\mathrm e} c)
 \bigl[-\ee^{+\ii p\cdot(x-y)/\hbar}\bigr]
 &=(\slashed p-m_{\mathrm e} c)\ee^{+\ii p\cdot(x-y)/\hbar}.
\end{align*}
把两式与 Pauli--Jordan 核的定义比较，即得到\eqnrefs{eq:dirac-PJ-factor-r8}。
\end{proof}


费米场的时序积在交换两个费米算符时必须带负号：
\begin{equation}
 \Tord\hat\psi(x)\bar{\hat\psi}(y)
 =\Theta(x^0-y^0)\hat\psi(x)\bar{\hat\psi}(y)
 -\Theta(y^0-x^0)\bar{\hat\psi}(y)\hat\psi(x).
 \label{eq:fermion-time-order-r2}
\end{equation}
由此定义 Dirac 的 Feynman 二点函数
\begin{equation*}
 S_F(x-y)
 \equiv\langle0|\Tord\hat\psi(x)\bar{\hat\psi}(y)|0\rangle.
\end{equation*}
对两个时间排序区域分别使用正、负频模，再由 $p^0$ 复平面上的 Feynman 边界条件合并，可得统一的四动量表示
\begin{equation}
 \boxed{
 S_F(x-y)
 =\int\frac{\dd^4 p}{(2\pi\hbar)^4}
 \frac{\ii\hbar(\slashed p+m_{\mathrm e} c)}{p^2-m_{\mathrm e}^2c^2+\ii0^+}
 \ee^{-\ii p\cdot(x-y)/\hbar}.}
 \label{eq:dirac-prop-struct-r2}
\end{equation}
其分子来自 Clifford 代数恒等式
\begin{equation*}
 (\slashed p-m_{\mathrm e} c)(\slashed p+m_{\mathrm e} c)=p^2-m_{\mathrm e}^2c^2,
\end{equation*}
所以传播子是自由一阶 Dirac 算符在 Feynman 边界条件下的 Green 核：
\begin{equation}
 \boxed{
 (\ii\hbar\slashed\partial_x-m_{\mathrm e} c)S_F(x-y)
 =\ii\hbar\,\delta^{(4)}(x-y)\id(4).}
 \label{eq:dirac-green-exact-r2}
\end{equation}

\begin{proposition}[Dirac Green 核的算符逆]
采用正文 Feynman 边界值处方与 Clifford 关系；时序微分的等时接触项由 CAR 固定。则有以下结论（原式 \eqnrefs{eq:dirac-green-exact-r2}）：
\begin{equation*}
 \boxed{
 (\ii\hbar\slashed\partial_x-m_{\mathrm e} c)S_F(x-y)
 =\ii\hbar\,\delta^{(4)}(x-y)\id(4).}
\end{equation*}
\end{proposition}
\begin{proof}
再计算时序二点函数的两个时间分支。真空条件 $b|0\rangle=d|0\rangle=0$ 使非零的 Wightman 核仅剩

\begin{align*}
 \langle0|\hat\psi(x)\bar{\hat\psi}(y)|0\rangle
 &=\int\frac{\dd^3 p}{(2\pi\hbar)^3}\frac{c}{2E_{\mathbf p}}
 (\slashed p+m_{\mathrm e} c)\ee^{-\ii p\cdot(x-y)/\hbar},\\
 \langle0|\bar{\hat\psi}(y)\hat\psi(x)|0\rangle
 &=\int\frac{\dd^3 p}{(2\pi\hbar)^3}\frac{c}{2E_{\mathbf p}}
 (\slashed p-m_{\mathrm e} c)\ee^{+\ii p\cdot(x-y)/\hbar}.
\end{align*}
时序乘积的第二项因为交换两个费米算符多出负号。把第一项解释为 $p^0=+E_{\mathbf p}/c$ 极点的留数，把第二项解释为 $p^0=-E_{\mathbf p}/c$ 极点的留数，并采用 Feynman 处方
\begin{equation*}
 p^2-m_{\mathrm e}^2c^2\longrightarrow p^2-m_{\mathrm e}^2c^2+\ii0^+,
\end{equation*}
两个三维质量壳积分合并为
\begin{equation*}
 S_F(x-y)
 =\int\frac{\dd^4 p}{(2\pi\hbar)^4}
 \frac{\ii\hbar(\slashed p+m_{\mathrm e} c)}{p^2-m_{\mathrm e}^2c^2+\ii0^+}
 \ee^{-\ii p\cdot(x-y)/\hbar},
\end{equation*}
即\eqnrefs{eq:dirac-prop-struct-r2}。

最后作用自由 Dirac 算符：
\begin{align*}
 (\ii\hbar\slashed\partial_x-m_{\mathrm e} c)S_F(x-y)
 &=\int\frac{\dd^4 p}{(2\pi\hbar)^4}\ii\hbar\,
 \frac{(\slashed p-m_{\mathrm e} c)(\slashed p+m_{\mathrm e} c)}{p^2-m_{\mathrm e}^2c^2+\ii0^+}
 \ee^{-\ii p\cdot(x-y)/\hbar}\\
 &=\ii\hbar\int\frac{\dd^4 p}{(2\pi\hbar)^4}
 \ee^{-\ii p\cdot(x-y)/\hbar}\\
 &=\ii\hbar\,\delta^{(4)}(x-y)\id(4),
\end{align*}
从而得到\eqnrefs{eq:dirac-green-exact-r2}。
\end{proof}

这说明传播子分子并非独立的记忆规则：同一个 $\slashed p+m_{\mathrm e} c$ 既由一阶矩阵算符的代数逆产生，也由自旋完备关系产生。

同一二点函数还可以从费米型高斯积分得到。由于采用 $x^0=ct$，作用量写成 $S=c^{-1}\int\dd^4 x\,\mathcal L$。定义具有能量量纲的自由 Dirac 算符
\begin{equation}
 \mathsf D_{0,x}\equiv\ii\hbar c\slashed\partial_x-m_{\mathrm e} c^2,
 \label{eq:dirac-energy-kernel-dp42}
\end{equation}
并以 Feynman 边界条件定义其右逆核
\begin{equation}
 \mathsf D_{0,x}\mathsf D_0^{-1}(x,y)
 =\delta^{(4)}(x-y)\id(4).
 \label{eq:dirac-energy-kernel-inverse-def-dp42}
\end{equation}
这两个 Green 核只差自由 Dirac 算符采用“动量量纲”还是“能量量纲”的归一化。二者的对应关系为
\begin{equation}
 \boxed{
 \mathsf D_0^{-1}(x,y)
 =\frac{S_F(x-y)}{\ii\hbar c}.}
 \label{eq:dirac-kernel-propagator-relation-dp42}
\end{equation}
这个归一化因子可直接由
$\mathsf D_{0,x}=c(\ii\hbar\slashed\partial_x-m_{\mathrm e} c)$ 与
\eqnrefs{eq:dirac-green-exact-r2}，
\begin{equation*}
 \mathsf D_{0,x}\!\left[\frac{S_F(x-y)}{\ii\hbar c}\right]
 =\delta^{(4)}(x-y)\id(4),
\end{equation*}
与右逆定义\eqnrefs{eq:dirac-energy-kernel-inverse-def-dp42} 比较即得到
\eqnrefs{eq:dirac-kernel-propagator-relation-dp42}。
场量纲满足 $[\psi]=[\bar\psi]=\mathrm{m^{-3/2}}$，而 $[\mathsf D_0^{-1}]=\mathrm{J^{-1}m^{-4}}$。若 Grassmann 外源 $\mathcal J_\psi$ 与 $\overline{\mathcal J}_\psi$ 满足 $[\mathcal J_\psi]=[\overline{\mathcal J}_\psi]=\mathrm{J\,m^{-3/2}}$，则自由生成泛函为
\begin{equation}
 \boxed{
 \frac{Z_0[\mathcal J_\psi,\overline{\mathcal J}_\psi]}{Z_0[0,0]}
 =\exp\!\left[
 -\frac{\ii}{\hbar c}
 \int\dd^4 x\,\dd^4 y\,
 \overline{\mathcal J}_\psi(x)\mathsf D_0^{-1}(x,y)\mathcal J_\psi(y)
 \right].}
 \label{eq:dirac-free-generating-functional-dp42}
\end{equation}

\begin{proposition}[Grassmann 高斯源的生成泛函]
设有限模 Grassmann 二次核可逆，采用正文源次序与左右微分约定；连续极限保持同一核和正则化。则有以下结论（原式 \eqnrefs{eq:dirac-free-generating-functional-dp42}）：
\begin{equation*}
 \boxed{
 \frac{Z_0[\mathcal J_\psi,\overline{\mathcal J}_\psi]}{Z_0[0,0]}
 =\exp\!\left[
 -\frac{\ii}{\hbar c}
 \int\dd^4 x\,\dd^4 y\,
 \overline{\mathcal J}_\psi(x)\mathsf D_0^{-1}(x,y)\mathcal J_\psi(y)
 \right].}
\end{equation*}
\end{proposition}
\begin{proof}
含 Grassmann 外源的自由作用量为

\begin{equation*}
 S_0[\psi,\bar\psi;\mathcal J_\psi,\overline{\mathcal J}_\psi]
 =\frac1c\int\dd^4 x\,
 \left[
 \bar\psi\mathsf D_0\psi+\overline{\mathcal J}_\psi\psi+\bar\psi\mathcal J_\psi
 \right].
\end{equation*}
相应泛函积分是
\begin{equation*}
 Z_0[\mathcal J_\psi,\overline{\mathcal J}_\psi]
 =\int\mathcal D\bar\psi\,\mathcal D\psi\,
 \exp\!\left\{\frac{\ii}{\hbar c}
 \int\dd^4 x\,
 [\bar\psi\mathsf D_0\psi+\overline{\mathcal J}_\psi\psi+\bar\psi\mathcal J_\psi]\right\}.
\end{equation*}
连续时空指标下的 Grassmann 高斯积分仍可作平移，取
\begin{equation*}
 \psi=\psi'-\mathsf D_0^{-1}\mathcal J_\psi,
 \qquad
 \bar\psi=\bar\psi'-\overline{\mathcal J}_\psi\mathsf D_0^{-1}.
\end{equation*}
逐项展开二次型：
\begin{align*}
 \bar\psi\mathsf D_0\psi
 +\overline{\mathcal J}_\psi\psi+\bar\psi\mathcal J_\psi
 ={}&
 (\bar\psi'-\overline{\mathcal J}_\psi\mathsf D_0^{-1})
 \mathsf D_0(\psi'-\mathsf D_0^{-1}\mathcal J_\psi)\\
 &+\overline{\mathcal J}_\psi(\psi'-\mathsf D_0^{-1}\mathcal J_\psi)
 +(\bar\psi'-\overline{\mathcal J}_\psi\mathsf D_0^{-1})\mathcal J_\psi\\
 ={}&\bar\psi'\mathsf D_0\psi'
 -\overline{\mathcal J}_\psi\mathsf D_0^{-1}\mathcal J_\psi.
\end{align*}
第二步使用 $\mathsf D_0\mathsf D_0^{-1}=1$，两个线性项分别相消。Grassmann 平移的雅可比行列式为 $1$，所以与外源无关的高斯积分等于 $Z_0[0,0]$，余下外源项给出\eqnrefs{eq:dirac-free-generating-functional-dp42}。

量纲检查为
\begin{equation*}
 [\mathcal J_\psi][\overline{\mathcal J}_\psi][\mathsf D_0^{-1}]\,\mathrm{m^8}
 =(\mathrm{J^2m^{-3}})(\mathrm{J^{-1}m^{-4}})\mathrm{m^8}
 =\mathrm{J\,m}
 =[\hbar c].
\end{equation*}
因此生成泛函指数无量纲；这也排除了把 $\mathsf D_0^{-1}$ 写成 $S_F/(\ii\hbar)$ 的可能性。
\end{proof}

因此正则量子化得到的时序二点函数、Grassmann 源泛函的二阶导数，以及自由 Dirac 二次核的 Feynman 逆，都是同一个 Green 核的不同表示。

\subsection{Coulomb 规范}

经典 Maxwell 方程只含两个传播偏振，但四势 \(A^\mu\) 却有四个分量；多出来的自由度不是新粒子，而是同一电磁场的规范冗余。量子化时如果把四个分量都当作独立谐振子，会同时破坏约束和物理态计数。Coulomb 规范的价值在于先把瞬时约束与两个横向传播自由度分开，使物理光子 Hilbert 空间可以直接看见。

Maxwell 场不能像三个独立标量场那样直接量子化，因为四势含有规范冗余。三维势满足
\begin{equation*}
\mathbf E=-\nabla\phi-\partial_t\mathbf A,
\qquad
\mathbf B=\nabla\times\mathbf A,
\end{equation*}
而 Maxwell 拉氏量对 $\dot\phi$ 完全不依赖，所以 $\Pi_\phi=0$ 是初级约束。对 $\phi$ 变分又给出 Gauss 定律；无源区因此只剩横向电磁自由度。Coulomb 规范
\begin{equation*}
\nabla\cdot\mathbf A=0
\end{equation*}
直接把这一物理相空间显式化。

约束意味着正则括号也必须投影到横向子空间。定义
\begin{equation*}
(P_T)_{ij}=\delta_{ij}-\frac{\partial_i\partial_j}{\nabla^2},
\end{equation*}
其中 $\nabla^{-2}$ 按给定边界条件定义。物理变量 $A^T=P_TA$、$\Pi^T=P_T\Pi$ 满足
\begin{equation}
\boxed{
[A_i(\mathbf x),\Pi_j(\mathbf y)]
=\ii\hbar\,\delta^T_{ij}(\mathbf x-\mathbf y),}
\label{eq:transverse-comm-r2}
\end{equation}
其中
\begin{equation*}
\delta^T_{ij}(\mathbf r)
=\int\frac{\dd^3 \kappa}{(2\pi)^3}
\left(\delta_{ij}-\hat\kappa_i\hat\kappa_j\right)\ee^{i\boldsymbol\kappa\cdot\mathbf r}.
\end{equation*}
因此光子的两个偏振并不是从 $A^\mu$ 中“手工删掉”两个分量，而是规范约束后横向物理相空间的两个独立方向。

对每个自由波矢 $\boldsymbol\kappa$ 选两个横向单位矢量 $\boldsymbol\epsilon_\lambda$，$\lambda=1,2$，满足 $\boldsymbol\kappa\cdot\boldsymbol\epsilon_\lambda=0$ 与 $\boldsymbol\epsilon_\lambda^*\cdot\boldsymbol\epsilon_{\lambda'}=\delta_{\lambda\lambda'}$。它们在横向平面上的完备关系为
\begin{equation*}
\sum_{\lambda=1}^2\epsilon_{\lambda i}\epsilon_{\lambda j}^*
=\delta_{ij}-\hat\kappa_i\hat\kappa_j.
\end{equation*}
于是矢势的正常模展开为
\begin{equation}
\boxed{
\hat{\mathbf A}(\mathbf r,t)
=\sum_{\lambda=1}^2\int\frac{\dd^3 \kappa}{(2\pi)^3}
\sqrt{\frac{\hbar}{2\varepsilon_0\omega_{\boldsymbol\kappa}}}
\left[
\boldsymbol\epsilon_\lambda(\boldsymbol\kappa)a_\lambda(\boldsymbol\kappa)
 \ee^{i\boldsymbol\kappa\cdot\mathbf r-\ii\omega_{\boldsymbol\kappa}t}
+\mathrm{h.c.}\right],}
\label{eq:maxwell-mode-r2}
\end{equation}
其中 $\omega_{\boldsymbol\kappa}=c|\boldsymbol\kappa|$，且
\begin{equation*}
[a_\lambda(\boldsymbol\kappa),a_{\lambda'}^\dagger(\boldsymbol\kappa')]
=(2\pi)^3\delta_{\lambda\lambda'}\delta^{(3)}(\boldsymbol\kappa-\boldsymbol\kappa').
\end{equation*}
哈密顿量因而是两个横向谐振子分支的直和：
\begin{equation*}
H_{\rm EM}
=\sum_{\lambda=1}^2\int\frac{\dd^3 \kappa}{(2\pi)^3}
\hbar\omega_{\boldsymbol\kappa}
\left(a_\lambda^\dagger a_\lambda+\frac12\right).
\end{equation*}
单模相干态 $a|\alpha\rangle=\alpha|\alpha\rangle$ 的 $\langle\mathbf A\rangle$ 满足经典 Maxwell 方程；当相对量子涨落随平均光子数增大而受控减小时，量子辐射场连续过渡到半经典外场描述。

连续平面波展开已经固定真空光子的两个偏振。在腔和纳米结构中还要回答两个实际问题：一个光子的场幅度如何由能量归一化决定，以及具有同一平均能量的光态为何产生不同交换统计。

\subsubsection{无损正常模的光子归一化}
\label{sec:feqo-quantized-field}
\begin{proposition}[无损正常模的每光子归一化]
在无损、无色散、封闭的非磁性腔中，设实对称正定介电张量为 $\boldsymbol\varepsilon_r(\mathbf r)$。取频率 $\omega_\lambda>0$ 的实驻波模 $\mathbf f_\lambda$ 满足
\begin{align}
 \nabla\times\nabla\times\mathbf f_\lambda
 &=\frac{\omega_\lambda^2}{c^2}\boldsymbol\varepsilon_r\mathbf f_\lambda,
 \qquad \nabla\cdot(\boldsymbol\varepsilon_r\mathbf f_\lambda)=0,\nonumber\\
 \int\dd^3r\,\mathbf f_\lambda\cdot\boldsymbol\varepsilon_r\mathbf f_{\lambda'}
 &=\delta_{\lambda\lambda'}.
 \label{eq:feqo-normal-mode-energy-normalization}
\end{align}
边界取理想导体条件或其他使旋度分部积分表面项为零的自伴条件。单模坐标 $\mathbf A=Q_\lambda\mathbf f_\lambda/\sqrt{\varepsilon_0}$ 给
\begin{equation}
 H_\lambda=\tfrac12(P_\lambda^2+\omega_\lambda^2Q_\lambda^2)
 =\hbar\omega_\lambda(a_\lambda^\dagger a_\lambda+\tfrac12),
 \label{eq:feqo-single-mode-hamiltonian}
\end{equation}
并且该模的正频电场是
\begin{equation}
 \hat{\mathbf E}^{(+)}_\lambda(\mathbf r,t)
 =\ii\sqrt{\frac{\hbar\omega_\lambda}{2\varepsilon_0}}\,
 \mathbf f_\lambda(\mathbf r)a_\lambda\ee^{-\ii\omega_\lambda t}.
 \label{eq:feqo-single-mode-field}
\end{equation}
这些表达式只在上述能量与自伴条件下适用；有色散时需加入材料能量，有损时需保留连续环境。
\end{proposition}

\begin{proof}
令 $\mathbf E=-\dot Q_\lambda\mathbf f_\lambda/\sqrt{\varepsilon_0}$、$\mathbf B=Q_\lambda\nabla\times\mathbf f_\lambda/\sqrt{\varepsilon_0}$。模方程乘 $\mathbf f_\lambda$ 并积分，边界项为零，得到
\begin{equation*}
 \int\dd^3r\,|\nabla\times\mathbf f_\lambda|^2
 =\frac{\omega_\lambda^2}{c^2}
 \int\dd^3r\,\mathbf f_\lambda\cdot\boldsymbol\varepsilon_r\mathbf f_\lambda
 =\frac{\omega_\lambda^2}{c^2}.
\end{equation*}
电磁能 $\frac12\int(\varepsilon_0\mathbf E\cdot\boldsymbol\varepsilon_r\mathbf E+\mu_0^{-1}|\mathbf B|^2)\dd^3r$ 因而为 $(\dot Q_\lambda^2+\omega_\lambda^2Q_\lambda^2)/2$，正则动量为 $P_\lambda=\dot Q_\lambda$。使用\eqnrefs{eq:sho-ladder-r3} 的反解
\begin{equation*}
 Q_\lambda=\sqrt{\frac{\hbar}{2\omega_\lambda}}(a_\lambda+a_\lambda^\dagger),
 \qquad P_\lambda=-\ii\sqrt{\frac{\hbar\omega_\lambda}{2}}(a_\lambda-a_\lambda^\dagger)
\end{equation*}
得到 Hamiltonian；$a_\lambda(t)=a_\lambda\ee^{-\ii\omega_\lambda t}$ 再给电场。由于 $[\mathbf f_\lambda]=\mathrm{m^{-3/2}}$，\eqnrefs{eq:feqo-single-mode-field} 的单位为 $\mathrm{V\,m^{-1}}$。
\end{proof}

\subsubsection{正常模与横向 Green 谱}
\label{sec:feqo-green-spectral}
\begin{proposition}[正常模与横向 Green 谱的对应]
对上述封闭系统，传播子空间中的 Green 核为
\begin{equation}
 \mathbf G_T^R(\mathbf r,\mathbf r';\omega)
 =c^2\sum_\lambda
 \frac{\mathbf f_\lambda(\mathbf r)\otimes\mathbf f_\lambda^*(\mathbf r')}
 {\omega_\lambda^2-(\omega+\ii0^+)^2}.
 \label{eq:feqo-green-mode-expansion}
\end{equation}
在 $\omega>0$，
\begin{equation}
 \operatorname{Im}_{H}\mathbf G_T^R
 =\frac{\pi c^2}{2}\sum_\lambda
 \frac{\mathbf f_\lambda\otimes\mathbf f_\lambda^*}{\omega_\lambda}
 \delta(\omega-\omega_\lambda).
 \label{eq:feqo-green-mode-positive-spectrum}
\end{equation}
量子场与经典 Green 张量使用同一组模式及能量归一化；完整 Green 张量还包含由 Gauss 约束确定的纵向部分。
\end{proposition}

\begin{proof}
在加权内积 $\langle\mathbf u,\mathbf v\rangle_\varepsilon=\int\mathbf u^*\cdot\boldsymbol\varepsilon_r\mathbf v\,\dd^3r$ 下展开任意物理横向场。Maxwell 算符作用于单个模给
\begin{equation*}
 \mathsf L_{\rm EM}(\omega)\mathbf f_\lambda
 =\frac{\omega_\lambda^2-\omega^2}{c^2}\boldsymbol\varepsilon_r\mathbf f_\lambda.
\end{equation*}
逐模求逆并用完备关系得到谱和。$\operatorname{Im}(x-\ii0^+)^{-1}=\pi\delta(x)$ 和正频根的 Jacobian $1/(2\omega_\lambda)$ 给第二式。

把电场模展开中的真空收缩 $\langle0|a_\lambda a_{\lambda'}^\dagger|0\rangle=\delta_{\lambda\lambda'}$ 与谱和比较，可得正频谱分量的真空相关
\begin{equation*}
 \langle0|\hat{\mathbf E}^{(+)}(\mathbf r,\omega)
 \hat{\mathbf E}^{(-)}(\mathbf r',\omega')|0\rangle
 =\frac{\hbar\mu_0\omega^2}{\pi}
 \operatorname{Im}_{H}\mathbf G_T^R(\mathbf r,\mathbf r';\omega)
 \delta(\omega-\omega').
\end{equation*}
离散腔中该式在频率分布意义下读取；它也是后面连续储库涨落公式的无损正常模极限。
\end{proof}

\begin{exbox}[周期盒、连续归一化与每体积的光子态密度]
在边长 $L$、体积 $V=L^3$ 的真空周期盒中，$\boldsymbol\kappa=2\pi\mathbf n/L$，每个非零波矢有两个横向偏振。离散模满足 $[a_{\boldsymbol\kappa\lambda},a_{\boldsymbol\kappa'\lambda'}^\dagger]=\delta_{\boldsymbol\kappa\boldsymbol\kappa'}\delta_{\lambda\lambda'}$。每个格点占据波矢体积 $(2\pi)^3/V$，因此
\begin{equation*}
 \sum_{\boldsymbol\kappa}\longrightarrow\frac{V}{(2\pi)^3}\int\dd^3\kappa,
 \qquad \delta_{\boldsymbol\kappa\boldsymbol\kappa'}
 \longrightarrow\frac{(2\pi)^3}{V}\delta^{(3)}(\boldsymbol\kappa-\boldsymbol\kappa').
\end{equation*}
与\eqnrefs{eq:maxwell-mode-r2} 相同的连续算符是 $a_\lambda(\boldsymbol\kappa)=\sqrt V\,a_{\boldsymbol\kappa\lambda}$，故其对易子为 $(2\pi)^3\delta^{(3)}$，离散场前的 $V^{-1/2}$ 正好转成连续积分系数。若改用没有 $(2\pi)^3$ 的 delta 对易子，必须同时缩放算符和场系数，不能只改积分测度。

真实单光子态由归一化波包构造：
\begin{equation}
 |1_F\rangle=\sum_{\lambda=1}^2\int\frac{\dd^3\kappa}{(2\pi)^3}
 F_\lambda(\boldsymbol\kappa)a_\lambda^\dagger(\boldsymbol\kappa)|0\rangle,
 \qquad \sum_\lambda\int\frac{\dd^3\kappa}{(2\pi)^3}|F_\lambda|^2=1.
 \label{eq:feqo-normalized-photon-wavepacket}
\end{equation}
对易关系给 $\langle1_F|1_F\rangle=1$，平均激发能为
$\sum_\lambda\int\hbar c|\boldsymbol\kappa|\,|F_\lambda|^2\dd^3\kappa/(2\pi)^3$；单个确定波矢的态则是不可归一化的广义本征态。

在频率区间 $[\omega,\omega+\dd\omega]$ 计数两个偏振的格点，得到每单位体积的态密度
\begin{equation}
 \rho_\gamma(\omega)\dd\omega
 =\frac{2}{(2\pi)^3}4\pi\kappa^2\dd\kappa
 =\frac{\omega^2}{\pi^2c^3}\dd\omega.
 \label{eq:feqo-vacuum-photon-density}
\end{equation}
真空 Green 张量的独立核验是
\begin{equation*}
 \frac{2\omega}{\pi c^2}\Tr\operatorname{Im}\mathbf G_0^R(\mathbf r,\mathbf r;\omega)
 =\frac{2\omega}{\pi c^2}\frac{3\omega}{6\pi c}
 =\rho_\gamma(\omega).
\end{equation*}
这里用了\eqnrefs{eq:feqo-vacuum-dipole-radiation} 附近得到的真空局域谱；它把模式计数、每光子能量和 Green 响应的归一化连在一起。
\end{exbox}

\subsubsection{相干态与单模热态的涨落}
\label{sec:feqo-coherent-thermal}
\begin{definition}[相干态与单模热态]
相干态是单位范数的湮灭算符本征态 $a|\alpha\rangle=\alpha|\alpha\rangle$，$\alpha\in\mathbb C$。单模热态是温度 $T>0$ 下的 Gibbs 密度算符 $\rho_{\rm th}=\ee^{-H_\lambda/(k_BT)}/\Tr\ee^{-H_\lambda/(k_BT)}$。二者在数态基中分别为
\begin{align}
 |\alpha\rangle&=\ee^{-|\alpha|^2/2}\sum_{n=0}^\infty\frac{\alpha^n}{\sqrt{n!}}|n\rangle,
 \label{eq:feqo-coherent-state}\\
 \rho_{\rm th}&=(1-x)\sum_{n=0}^\infty x^n|n\rangle\langle n|,
 \qquad x=\ee^{-\hbar\omega_\lambda/(k_BT)},\quad
 \bar n=\frac{x}{1-x}.
 \label{eq:feqo-thermal-state}
\end{align}
相干态的平均振幅非零，热态则 $\langle a\rangle=0$，即使二者有相同 $\bar n$。
\end{definition}

对单模相干态，\eqnrefs{eq:feqo-single-mode-field} 给平均正频场幅度
$\mathbf e_+=\ii\sqrt{\hbar\omega_\lambda/(2\varepsilon_0)}\,\mathbf f_\lambda\alpha$。
完整实场是 $\mathbf e_+\ee^{-\ii\omega_\lambda t}+\mathbf e_+^*\ee^{\ii\omega_\lambda t}$。
若写成 $\operatorname{Re}[\boldsymbol{\mathcal E}\ee^{-\ii\omega_\lambda t}]$，则 $\boldsymbol{\mathcal E}=2\mathbf e_+$；正频幅度与实场复幅度之间的这个二倍关系同时影响电子耦合和时间平均功率。

\begin{proposition}[数涨落、正交分量与二阶相关]
令 $N=a^\dagger a$、$X=(a+a^\dagger)/\sqrt2$、$Y=(a-a^\dagger)/(\ii\sqrt2)$，则 $[X,Y]=\ii$。相干态满足
\begin{equation}
 \langle N\rangle=|\alpha|^2,\qquad
 \operatorname{Var}N=|\alpha|^2,\qquad
 \operatorname{Var}X=\operatorname{Var}Y=\tfrac12.
 \label{eq:feqo-coherent-variances}
\end{equation}
热态满足
\begin{equation}
 \operatorname{Var}N=\bar n(1+\bar n),\qquad
 \operatorname{Var}X=\operatorname{Var}Y=\bar n+\tfrac12.
 \label{eq:feqo-thermal-variances}
\end{equation}
对 $\bar n>0$，零延迟正规序相关 $g^{(2)}(0)=\langle N(N-1)\rangle/\bar n^2$ 依次是相干态的 $1$、热态的 $2$ 和数态 $|n\rangle$ 的 $1-1/n$。
\end{proposition}

\begin{proof}
相干态递推 $c_{n+1}\sqrt{n+1}=\alpha c_n$ 与范数固定 $c_0=\ee^{-|\alpha|^2/2}$，因此数概率是 Poisson 权重。对其产生函数
\begin{equation*}
 F_\alpha(z)=\langle z^N\rangle=\ee^{|\alpha|^2(z-1)}
\end{equation*}
在 $z=1$ 求一、二次导数，得 $\langle N\rangle=|\alpha|^2$ 和 $\langle N(N-1)\rangle=|\alpha|^4$。

热态中零点能在配分函数归一化时相消，故几何权重的产生函数为 $F_T(z)=(1-x)/(1-xz)$。其导数给 $\langle N\rangle=x/(1-x)$、$\langle N(N-1)\rangle=2x^2/(1-x)^2$。由 $N^2=N(N-1)+N$ 得两种方差。

相干态中 $\langle a\rangle=\alpha$、$\langle a^2\rangle=\alpha^2$、$\langle aa^\dagger\rangle=|\alpha|^2+1$。逐项展开 $X^2,Y^2$ 并减去均值平方，剩下各自 $1/2$。因此 $\Delta X\Delta Y=1/2$，或 $\Delta Q\Delta P=\hbar/2$。热态的数态对角性使 $\langle a\rangle=\langle a^2\rangle=0$，而 $\langle aa^\dagger\rangle=\bar n+1$，于是两个正交分量方差均为 $\bar n+1/2$。
\end{proof}

\begin{exbox}[同为四个光子的平均能量，交换统计仍不同]
比较 $|4\rangle$、$\alpha=2\ee^{\ii\theta}$ 的相干态，以及 $x=4/5$ 的热态。它们均有 $\langle N\rangle=4$，但数方差依次为 $0,4,20$，$\langle N(N-1)\rangle$ 依次为 $12,16,32$。

若在受控微扰区间内一光子吸收振幅正比于 $a$，两光子吸收振幅正比于 $a^2$，且其他电子核和失谐相同，那么一光子概率采样 $\langle a^\dagger a\rangle=4$，两光子概率则采样 $\langle a^{\dagger2}a^2\rangle=12,16,32$。所以仅匹配平均光子数无法匹配所有电子边带概率；进入强耦合时必须保留完整联合幺正演化。
\end{exbox}

\begin{definition}[电子电流与量子辐射场的耦合]
在 Coulomb 规范中，横向辐射的线性耦合为
\begin{equation}
 H_{I,T}(t)=-\int\dd^3r\,
 \hat{\mathbf j}_{\rm ch,T}(\mathbf r,t)\cdot\hat{\mathbf A}_T(\mathbf r,t).
 \label{eq:feqo-current-field-coupling}
\end{equation}
Gauss 约束同时产生纵向 Coulomb 能量，不能在具有纵向电荷响应的过程中舍去。单电子电荷流矩阵元是 $q_{\mathrm e}c\Psi_f^\dagger\boldsymbol\alpha\Psi_i$，与\eqnrefs{eq:qed-interaction-lagrangian} 的 QED 顶角使用相同的 Dirac 双线性。
\end{definition}


\subsection{协变规范量子化}

Coulomb 规范的优点是物理自由度一眼可见：只量子化两个横向模式。它的代价是 Lorentz 协变性不显式。协变散射论反过来选择“先保留四势的四个分量，再从扩大后的状态空间中筛出物理态”。因此这一小节真正需要回答两个问题：为什么扩大后的空间会出现负范数方向，以及规范条件怎样把这些方向从可观测物理中消掉。答案分别由不定 Minkowski 度规和 Gupta--Bleuler 条件给出。

协变散射论改用规范固定后的 Maxwell 拉氏量
\begin{equation}
 \mathcal L_{\rm EM,\xi_{\rm gf}}
 =-\frac1{4\mu_0}F_{\mu\nu}F^{\mu\nu}
 -\frac1{2\xi_{\rm gf}\mu_0}(\partial_\mu A^\mu)^2,
 \label{eq:cov-gauge-lagrangian-r2}
\end{equation}
其中无量纲参数 $\xi_{\rm gf}$ 只标记规范固定方式；$\xi_{\rm gf}=1$ 为 Feynman 规范。协变形式先把四势的四个分量都放入扩大后的状态空间，再用约束和零范数等价关系提取两个物理光子自由度。

在 Feynman 规范中，自由方程为 $\Box A^\mu=0$。对四波矢 $\kappa^\mu=(\omega_{\boldsymbol\kappa}/c,\boldsymbol\kappa)$ 选一组 Minkowski 正交基 $\epsilon_\lambda^\mu$，$\lambda=0,1,2,3$，并取
\begin{equation}
 \epsilon_\lambda\cdot\epsilon_{\lambda'}^*=\eta_{\lambda\lambda'}.
 \label{eq:fourpol-orthonormal-r13}
\end{equation}
四势可展开为
\begin{equation}
 A^\mu(x)=
 \sum_{\lambda=0}^{3}
 \int\frac{\dd^3 \kappa}{(2\pi)^3}
 \sqrt{\frac{\hbar}{2\varepsilon_0\omega_{\boldsymbol\kappa}}}
 \left[
 \epsilon_\lambda^\mu a_\lambda \ee^{-\ii\kappa\cdot x}
 +\epsilon_\lambda^{\mu *}a_\lambda^\dagger \ee^{+\ii\kappa\cdot x}
 \right].
 \label{eq:covariant-photon-fourpol-r7}
\end{equation}
省略的算符自变量均为 $\boldsymbol\kappa$。基的完备性给
\begin{equation}
 \boxed{
 \sum_{\lambda,\lambda'}
 \eta^{\lambda\lambda'}
 \epsilon_\lambda^\mu\epsilon_{\lambda'}^{\nu *}
 =\eta^{\mu\nu}.}
 \label{eq:fourpol-completeness-r7}
\end{equation}
正则量子化相应要求
\begin{equation}
 \boxed{
 [a_\lambda(\boldsymbol\kappa),a_{\lambda'}^\dagger(\boldsymbol\kappa')]
 =-\eta_{\lambda\lambda'}(2\pi)^3
 \delta^{(3)}(\boldsymbol\kappa-\boldsymbol\kappa').}
 \label{eq:covariant-photon-mode-comm-r7}
\end{equation}
时间样模因而具有负范数，
\begin{equation}
 \langle0|a_0(\boldsymbol\kappa)a_0^\dagger(\boldsymbol\kappa')|0\rangle
 =-(2\pi)^3\delta^{(3)}(\boldsymbol\kappa-\boldsymbol\kappa'),
 \label{eq:timelike-negative-norm-r7}
\end{equation}
说明扩大后的协变态空间本身不是物理 Hilbert 空间。

把 $A^\mu=A^{\mu(+)}+A^{\mu(-)}$ 分成正、负频部分后，Gupta--Bleuler 条件选择物理态：
\begin{equation}
 \boxed{
 \partial_\mu A^{\mu(+)}(x)|\Psi_{\rm phys}\rangle=0.}
 \label{eq:gupta-bleuler-r2}
\end{equation}
若光子沿 $z$ 方向传播，并取 $\epsilon_{(0)}^\mu=(1,0,0,0)$、$\epsilon_{(3)}^\mu=(0,0,0,1)$，该条件化为
\begin{equation}
 \boxed{(a_3-a_0)|\Psi_{\rm phys}\rangle=0.}
 \label{eq:GB-mode-condition-r7}
\end{equation}
相应组合
\begin{equation*}
 |N\rangle=(a_3^\dagger-a_0^\dagger)|0\rangle
\end{equation*}
具有
\begin{equation}
 \boxed{\langle N|N\rangle=0,}
 \label{eq:GB-null-state-r7}
\end{equation}
所以时间样与纵向激发只产生零范数规范方向。对这些方向取商以后，物理 Hilbert 空间只留下两个横向光子偏振。

规范固定同时使 Maxwell 二次核可逆。对\eqnrefs{eq:cov-gauge-lagrangian-r2} 分部积分可写成
\begin{equation}
 S_{\rm EM,\xi_{\rm gf}}^{(2)}
 =\frac1{2c}\int\dd^4 x\,
 A_\mu\frac1{\mu_0}
 \left[
 \eta^{\mu\nu}\Box
 -\left(1-\frac1{\xi_{\rm gf}}\right)
 \partial^\mu\partial^\nu
 \right]A_\nu.
 \label{eq:maxwell-gf-quadratic-r6}
\end{equation}
在四波矢空间定义纵向和横向投影算符
\begin{equation*}
 P_L^{\mu\nu}(\kappa)=\frac{\kappa^\mu\kappa^\nu}{\kappa^2},
 \qquad
 P_T^{\mu\nu}(\kappa)=\eta^{\mu\nu}-P_L^{\mu\nu}(\kappa),
\end{equation*}
则二次核为
\begin{equation}
 \boxed{
 \mathcal K_{\xi_{\rm gf}}^{\mu\nu}(\kappa)
 =-\frac{\kappa^2}{\mu_0}
 \left(P_T^{\mu\nu}+\frac1{\xi_{\rm gf}}P_L^{\mu\nu}\right).}
 \label{eq:maxwell-kernel-projectors-r6}
\end{equation}
投影代数
\begin{equation}
 P_T^2=P_T,
 \qquad
 P_L^2=P_L,
 \qquad
 P_TP_L=P_LP_T=0,
 \qquad
 P_T+P_L=I
 \label{eq:maxwell-projector-algebra-r13}
\end{equation}
把四维矢量空间分成两个互补本征子空间，因此核的逆为
\begin{equation}
 \boxed{
 (\mathcal K_{\xi_{\rm gf}}^{-1})^{\mu\nu}(\kappa)
 =-\frac{\mu_0}{\kappa^2}
 \left(P_T^{\mu\nu}+\xi_{\rm gf}P_L^{\mu\nu}\right).}
 \label{eq:maxwell-kernel-inverse-r6}
\end{equation}
施加 Feynman 边界值后，真空时序二点函数为
\begin{equation}
 \boxed{
 \widetilde D_F^{\mu\nu}(\kappa)
 =\frac{-\ii\mu_0\hbar}{\kappa^2+\ii0^+}
 \left[
 \eta^{\mu\nu}-(1-\xi_{\rm gf})
 \frac{\kappa^\mu\kappa^\nu}{\kappa^2+\ii0^+}
 \right].}
 \label{eq:photon-prop-xi-r2}
\end{equation}
Feynman 规范把它化为
\begin{equation}
 \boxed{
 \widetilde D_F^{\mu\nu}(\kappa)
 =\frac{-\ii\mu_0\hbar\eta^{\mu\nu}}
 {\kappa^2+\ii0^+}.}
 \label{eq:photon-prop-feynman-r2}
\end{equation}
物理四动量 $q^\mu=\hbar\kappa^\mu$ 与四波矢表示之间满足
\begin{equation}
 \boxed{
 \begin{aligned}
 D_F^{\mu\nu}(x-y)
 &=\int\frac{\dd^4 \kappa}{(2\pi)^4}
 \frac{-\ii\mu_0\hbar\eta^{\mu\nu}}
 {\kappa^2+\ii0^+}\ee^{-\ii\kappa\cdot(x-y)},\\
 &=\int\frac{\dd^4 q}{(2\pi\hbar)^4}
 \frac{-\ii\mu_0\hbar^3\eta^{\mu\nu}}
 {q^2+\ii0^+}\ee^{-\ii q\cdot(x-y)/\hbar}.
 \end{aligned}}
 \label{eq:photon-prop-kappa-to-q-r13}
\end{equation}
因此以物理四动量表示内部光子线时，传播子分子携带 $\hbar^3$。守恒流满足 $q_\mu j^\mu=0$，所以规范依赖的纵向项在规范不变矩阵元中消失。

协变路径积分还必须处理同一规范轨道被重复积分的问题。前一章已经在一般规范群上由规范轨道方向的泛函 Jacobian 建立了 Faddeev--Popov 恒等式；QED 这里只代入阿贝尔 $U(1)$ 的规范变换。取
\begin{equation*}
 G[A]=\partial_\mu A^\mu,
 \qquad
 A_\mu^\alpha=A_\mu-\partial_\mu\alpha,
\end{equation*}
则规范条件沿轨道的线性响应为
\begin{equation*}
 \frac{\delta G[A^\alpha](x)}{\delta\alpha(y)}
 =-\Box_x\delta^{(4)}(x-y).
\end{equation*}
因此
\begin{equation}
 \boxed{
 \Delta_{\rm FP}[A]=\det(-\Box),}
 \label{eq:FP-QED-det-r8}
\end{equation}
并且这个行列式与 $A_\mu$ 无关。于是阿贝尔鬼场只给出可吸收入归一化的自由二次行列式，不存在鬼场--光子顶角；非阿贝尔理论中轨道 Jacobian 含规范联络，鬼场相互作用才会出现。

把规范切片值写成 $G[A]=f$，再以宽度 $\xi_{\rm gf}$ 的 Gaussian 权重对辅助函数 $f(x)$ 积分，就得到同一规范轨道上的协变规范族：
若把与 $A_\mu$ 无关的泛函测度归一化常数记为 $\mathcal N_{\rm gf}$，则
\begin{equation}
 \boxed{
 \begin{aligned}
 &\int\mathcal Df\,
 \exp\!\left[-\frac{\ii}{2\hbar c\,\mu_0\xi_{\rm gf}}
 \int\dd^4 x\,f^2(x)\right]
 \delta\!\left(\partial_\mu A^\mu-f\right)\\
 &\qquad=
 \mathcal N_{\rm gf}
 \exp\!\left[-\frac{\ii}{2\hbar c\,\mu_0\xi_{\rm gf}}
 \int\dd^4 x\,(\partial_\mu A^\mu)^2\right].
 \end{aligned}}
 \label{eq:FP-gaussian-average-r8}
\end{equation}
\eqnrefs{eq:FP-gaussian-average-r8} 正好产生\eqnrefs{eq:cov-gauge-lagrangian-r2} 的规范固定项。$\xi_{\rm gf}$ 因而是规范切片族的参数；规范不变可观测量对它的依赖由 Ward/BRST 恒等式消去。

% -----------------------------------------------------------------------------
% Physical photon polarization projector after Maxwell quantization
% -----------------------------------------------------------------------------

Maxwell 约束把实光子的物理自由度限制为两个横向偏振。散射振幅中还需要把这两个偏振的求和写成协变张量。设实光子的物理四动量为 $k^\mu$，满足 $k^2=0$；再选一个满足 $k\!\cdot n\neq0$ 的辅助四矢量 $n^\mu$ 来固定偏振代表元。取
\begin{equation}
 k\!\cdot\epsilon_\lambda=0,
 \qquad
 n\!\cdot\epsilon_\lambda=0,
 \qquad
 \epsilon_\lambda\!\cdot\epsilon_{\lambda'}^*=-\delta_{\lambda\lambda'},
 \qquad \lambda=1,2.
 \label{eq:physical-photon-polarization-constraints-r68}
\end{equation}
两个物理偏振的完备张量为


\begin{proposition}[物理极化空间的协变完备投影]
设零四动量与所选辅助矢量满足正文条件，两个物理极化向量正交归一；相等在横向空间或规范等价类中理解。则有
\begin{equation}
 {
 \sum_{\lambda=1}^{2}
 \epsilon_\lambda^\mu(k)\epsilon_\lambda^{\nu *}(k)
 =-\eta^{\mu\nu}
 +\frac{k^\mu n^\nu+n^\mu k^\nu}{k\!\cdot n}
 -\frac{n^2k^\mu k^\nu}{(k\!\cdot n)^2}.}
\label{eq:physical-photon-polarization-projector-dp42}
\end{equation}
\end{proposition}
\begin{proof}
取实光子的物理四动量 $k^\mu$，满足 $k^2=0$，并选一个不与 $k$ 正交的辅助四矢量 $n^\mu$。物理偏振 $\epsilon_\lambda^\mu$ 满足

\begin{equation*}
 k\!\cdot\epsilon_\lambda=0,
 \qquad
 n\!\cdot\epsilon_\lambda=0,
 \qquad
 \epsilon_\lambda\!\cdot\epsilon_{\lambda'}^*=-\delta_{\lambda\lambda'},
 \qquad \lambda=1,2.
\end{equation*}
偏振和
\begin{equation*}
 P^{\mu\nu}(k,n)
 \equiv\sum_{\lambda=1,2}
 \epsilon_\lambda^\mu\epsilon_\lambda^{\nu*}
\end{equation*}
必须由 $\eta^{\mu\nu},k^\mu,n^\mu$ 构成，并满足
\begin{equation*}
 P^{\mu\nu}k_\nu=0,
 \qquad
 P^{\mu\nu}n_\nu=0.
\end{equation*}
从一般张量设式
\begin{equation*}
 P^{\mu\nu}
 =-\eta^{\mu\nu}
 +A(k^\mu n^\nu+n^\mu k^\nu)
 +B k^\mu k^\nu
 +C n^\mu n^\nu
\end{equation*}
出发。与 $k_\nu$ 收缩并用 $k^2=0$：
\begin{equation*}
 0=-k^\mu+A(k\!\cdot n)k^\mu+C(k\!\cdot n)n^\mu.
\end{equation*}
由于 $k^\mu,n^\mu$ 独立，
\begin{equation*}
 A=\frac1{k\!\cdot n},
 \qquad C=0.
\end{equation*}
再与 $n_\nu$ 收缩：
\begin{equation*}
 0=-n^\mu+\frac{k^\mu n^2+n^\mu(k\!\cdot n)}{k\!\cdot n}
 +B(k\!\cdot n)k^\mu,
\end{equation*}
从 $k^\mu$ 系数得到
\begin{equation*}
 B=-\frac{n^2}{(k\!\cdot n)^2}.
\end{equation*}
因此
\begin{equation*}
 {
 P^{\mu\nu}(k,n)
 =-\eta^{\mu\nu}
 +\frac{k^\mu n^\nu+n^\mu k^\nu}{k\!\cdot n}
 -\frac{n^2k^\mu k^\nu}{(k\!\cdot n)^2}.}
 \end{equation*}
若振幅的光子指标与守恒流 $J_\mu$ 收缩，并满足 $k^\mu J_\mu=0$，则
\begin{equation*}
 {
 J_\mu P^{\mu\nu}(k,n)J_\nu^*
 =-J_\mu J^{\mu*}.}
 \end{equation*}
所有依赖辅助矢量 $n^\mu$ 的项都消失。因此在规范不变可观测量中，物理偏振求和常可在与守恒振幅收缩之后简写为 $-\eta^{\mu\nu}$；这不是偏振空间本身的恒等式。
\end{proof}

辅助矢量 $n^\mu$ 只选择规范代表元，不进入规范不变的可观测量。若光子指标与守恒流 $J_\mu$ 收缩，$k^\mu J_\mu=0$ 使所有含 $k^\mu$ 的项消失，因此
\begin{equation}
 {
 J_\mu
 \left(\sum_{\lambda=1}^{2}
 \epsilon_\lambda^\mu\epsilon_\lambda^{\nu *}\right)
 J_\nu^*
 =-J_\mu J^{\mu *}.}
 \label{eq:physical-photon-polarization-conserved-current-dp42}
\end{equation}
所以常用的 $-\eta^{\mu\nu}$ 替换只在与守恒振幅收缩以后成立；它不是两个物理偏振子空间在完整四维 Minkowski 空间中的恒等算符。

\begin{exbox}[$z$ 方向光子的偏振与螺旋度]
取
\begin{equation*}
 k^\mu=\left(\frac{\hbar\omega}{c},0,0,\frac{\hbar\omega}{c}\right),
 \qquad k^2=0,
\end{equation*}
并在辐射规范代表元中取两个线偏振
\begin{equation*}
 \epsilon_x^\mu=(0,1,0,0),
 \qquad
 \epsilon_y^\mu=(0,0,1,0).
\end{equation*}
检验偏振的横向性与完备关系，再构造圆偏振基并求绕传播轴转动的本征值。
\textbf{解。} 逐项可检验
\begin{equation*}
 k\!\cdot\epsilon_x=k\!\cdot\epsilon_y=0,
 \qquad
 \epsilon_x^2=\epsilon_y^2=-1,
 \qquad
 \epsilon_x\!\cdot\epsilon_y=0.
\end{equation*}
因此这两个四矢量确实张成与传播方向垂直的二维物理偏振空间。它们的显式完备和为
\begin{equation*}
 \epsilon_x^\mu\epsilon_x^\nu
 +\epsilon_y^\mu\epsilon_y^\nu
 =\operatorname{diag}(0,1,1,0),
\end{equation*}
这正是\eqnrefs{eq:physical-photon-polarization-projector-dp42} 在 $n^\mu=(1,0,0,0)$ 下的结果。

对角化绕传播轴的转动，更自然的基是圆偏振
\begin{equation}
 \epsilon_\lambda^\mu
 =\frac1{\sqrt2}
 \left(0,1,\ii\lambda,0\right),
 \qquad \lambda=\pm1.
 \label{eq:photon-circular-polarizations-dp62}
\end{equation}
三维矢量表示的 $z$ 轴转动生成元满足
\begin{equation*}
 (S_z)_{ij}=-\ii\epsilon_{3ij},
 \qquad
 S_z\,\vb{\epsilon}_\lambda
 =\lambda\vb{\epsilon}_\lambda.
\end{equation*}
所以有限转动给出
\begin{equation*}
 \ee^{-\ii\phi S_z}\vb{\epsilon}_\lambda
 =\ee^{-\ii\lambda\phi}\vb{\epsilon}_\lambda.
\end{equation*}
这说明 $\lambda=\pm1$ 正是沿 $\vb{k}$ 的两个光子螺旋度。线偏振与圆偏振并不是两类不同光子，而只是同一个二维物理偏振空间的两组基：
\begin{equation*}
 \epsilon_x=\frac{\epsilon_++\epsilon_-}{\sqrt2},
 \qquad
 \epsilon_y=\frac{\epsilon_+-\epsilon_-}{\ii\sqrt2}.
\end{equation*}
散射实验若不分辨光子偏振，求和的是这个二维空间的一组完备基，因此用线偏振基或螺旋度基计算必须给出相同的未偏振截面。
\end{exbox}


\begin{exbox}[有质量矢量场为什么多一个纵向偏振]
为了和光子比较，考虑质量为 $M>0$ 的自由 Proca 模，取四动量
\begin{equation*}
 k^\mu=\left(\frac Ec,0,0,p\right),
 \qquad
 k^2=M^2c^2,
 \qquad
 E^2=M^2c^4+c^2p^2.
\end{equation*}
求第三个物理偏振及三态完备和，检验固定非零动量下的无质量极限。
\textbf{解。} 两个横向偏振仍可取 $\epsilon_x^\mu$、$\epsilon_y^\mu$。但此时还存在第三个满足 $k\cdot\epsilon_L=0$ 的归一化纵向偏振：
\begin{equation}
 {
 \epsilon_L^\mu(k)
 =\left(\frac{p}{Mc},0,0,\frac{E}{Mc^2}\right).}
 \label{eq:massive-vector-longitudinal-polarization-dp62}
\end{equation}
逐项检查得到
\begin{align*}
 k\cdot\epsilon_L
 &=\frac Ec\frac{p}{Mc}-p\frac{E}{Mc^2}=0,\\
 \epsilon_L^2
 &=\frac{p^2}{M^2c^2}-\frac{E^2}{M^2c^4}=-1.
\end{align*}
因此有质量自旋 $1$ 粒子具有三个物理偏振。三者的完备和为
\begin{equation}
 {
 \sum_{\lambda=1,2,L}
 \epsilon_\lambda^\mu(k)\epsilon_\lambda^{\nu *}(k)
 =-\eta^{\mu\nu}+\frac{k^\mu k^\nu}{M^2c^2}.}
 \label{eq:massive-vector-polarization-sum-dp62}
\end{equation}
由上式的显式分量 $\epsilon_L^0=p/(Mc)$ 与 $\vb{\epsilon}_L=(E/Mc^2)\hat{\vb{p}}$ 可见，在固定非零三动量 $p$ 下取 $M\to0$ 时，各分量均为 $O(M^{-1})$，因而有质量矢量场的纵向偏振不存在逐分量平滑的无质量极限。无质量规范场出现规范冗余；规范约束把纵向和时间样自由度从物理 Hilbert 空间中移除，只留下前一个算例中的两个螺旋度。标准模型的 $W^\pm,Z$ 因 Higgs 机制有质量，纵向偏振成为真实自由度；它与 Higgs 场中沿真空对称轨道的切向标量涨落之间的关系由破缺规范理论的质量矩阵和约束结构确定。
\end{exbox}

% LSZ 条件、约化公式和独立证明在正文中相邻排列。
实验真正制备和探测的是进入相互作用区之前、离开相互作用区之后能够独立传播的粒子；而场论计算首先给出的却是若干时空点上的场关联函数。两者之间因此还缺少一步：必须从相关函数中取出与一个真实外部粒子对应的传播部分，只留下相互作用本身的振幅。

对稳定粒子，这一步可以从精确二点函数的质量壳结构看清楚。当四动量接近真实粒子的质量壳时，二点函数出现一个单粒子极点；极点的位置给出物理质量，极点的留数给出场算符与规范化单粒子态之间的重叠。LSZ 约化所做的就是用自由逆运动算符消去这段外部传播，再按极点留数归一化，并把结果投影到真实电子自旋态或真实光子的横向偏振态。得到的剩余部分才是散射振幅。

因此，“截肢外腿”不是在图上把最外面的一条线擦掉，而是对质量壳极点执行一个明确的代数留数运算。内部传播线并不代表渐近探测到的单粒子态，所以不进行同样的质量壳投影。


光子外腿使用物理四动量 $k^\mu=\hbar\kappa^\mu$。把\eqnrefs{eq:maxwell-kernel-projectors-r6} 的四波矢核改写成物理四动量核，定义
\begin{equation}
 \mathscr K_{\xi_{\rm gf}}^{\mu\nu}(k)
 \equiv
 \mathcal K_{\xi_{\rm gf}}^{\mu\nu}(\kappa=k/\hbar)
 =-\frac{k^2}{\mu_0\hbar^2}
 \left(P_T^{\mu\nu}+\frac1{\xi_{\rm gf}}P_L^{\mu\nu}\right).
 \label{eq:photon-physical-momentum-inverse-kernel-dp42}
\end{equation}
由\eqnrefs{eq:photon-prop-kappa-to-q-r13}，自由核与物理四动量表示的二点函数满足 $\mathscr K_{\mu\rho}\widetilde D_F^{\rho\nu}=\ii\hbar\,\delta_\mu{}^\nu$。电子传播子同样满足自由 Dirac 逆核乘传播子等于 $\ii\hbar$。因此两类外腿的截肢算符都必须含 $(\ii\hbar)^{-1}$：
\begin{equation}
 \boxed{
 \begin{aligned}
 \text{电子外腿:}&\quad
 \widetilde S(p)\ \xrightarrow{\;[Z_{\rm pole}^{(e)}]^{-1/2}(\slashed p-mc)/(\ii\hbar)\;}\
 u,\bar u,v,\bar v,\\
 \text{光子外腿:}&\quad
 \widetilde D_{\mu\nu}(k)\
 \xrightarrow{\;[Z_{\rm pole}^{(\gamma)}]^{-1/2}\mathscr K^{\mu\rho}(k)/(\ii\hbar)\;}\
 \varepsilon_\rho^{(\lambda)}(k),
 \qquad k\cdot\varepsilon^{(\lambda)}=0.
 \end{aligned}}
 \label{eq:lsz-dirac-photon-mainline-dp8}
\end{equation}
这里用 $Z_{\rm pole}^{(e)}$ 与 $Z_{\rm pole}^{(\gamma)}$ 专门表示当前所用场变量的电子与横向光子单粒子极点权重。它们是二点函数的谱数据，负责把局域场算符产生的状态归一化成单位范数渐近粒子。重整化中常用的 $Z_2,Z_3$ 则表示“裸场与重整化场之间的场强重整化常数”；在一般方案中，它们与单粒子极点权重是不同对象。若采用在壳方案，$Z_2,Z_3$ 被选择来使重整化传播子的相应物理极点留数等于 $1$。

\begin{figure}[htbp]
 \centering
 \begin{tikzfig}[x=1cm,y=1cm,>=Stealth]
 \node[font=\bfseries] at (-3.5,1.55) {(a) full Green function};
 \node[font=\bfseries] at (3.55,1.55) {(b) amputated pole};

 % electron row
 \draw[fermion] (-5.3,.55)--(-2.7,.55);
 \node[above] at (-4.0,.55) {$S(p)$};
 \node[draw,rounded corners=3pt,minimum width=1.45cm,minimum height=.9cm] at (-1.95,.55) {$\Gamma$};
 \node at (0,.55) {$\xrightarrow{\quad Z_2^{-1/2}(\slashed p-mc)\quad}$};
 \node at (2.15,.55) {$\bar u_s(p)$};
 \draw[fermion] (2.75,.55)--(3.35,.55);
 \node[draw,rounded corners=3pt,minimum width=1.45cm,minimum height=.9cm] at (4.1,.55) {$\Gamma$};

 % photon row
 \draw[photon] (-5.3,-1.0)--(-2.7,-1.0);
 \node[above] at (-4.0,-1.0) {$D_{\mu\nu}(k)$};
 \node[draw,rounded corners=3pt,minimum width=1.45cm,minimum height=.9cm] at (-1.95,-1.0) {$\Gamma^\nu$};
 \node at (0,-1.0) {$\xrightarrow{\quad Z_3^{-1/2}\mathcal K^{\rho\mu}_\xi\quad}$};
 \node at (2.05,-1.0) {$\varepsilon_\rho^{(\lambda)*}$};
 \draw[photon] (2.85,-1.0)--(3.35,-1.0);
 \node[draw,rounded corners=3pt,minimum width=1.45cm,minimum height=.9cm] at (4.1,-1.0) {$\Gamma^\rho$};
\end{tikzfig}

 \caption{LSZ 约化的图形表示。完整 Green 函数含外腿单粒子极点；自由逆运动核除以 $\ii\hbar$ 后消去该传播极点，再投影到壳上电子自旋子或横向光子偏振。$\Gamma$ 表示其余连通截肢核。}
 \label{fig:qed-lsz-amputation-dp8}
\end{figure}

\begin{proposition}[稳定 Dirac 单粒子极点的 LSZ 约化]
设精确二点函数有稳定、孤立的单粒子极点，$Z_{\rm pole}^{(e)}>0$。
连通 Green 函数满足 $\widetilde G=\widetilde S\widetilde\Gamma+\widetilde G_{\rm reg}$，原始截肢核与余项在极点邻域解析。
以下 $m=m_{\rm phys}$；限式先提取极点，再与对应质量壳动量的旋量收缩，
$u_s(p)$ 的记号表示壳上投影，并不假设离壳 $p$ 也有自由 Dirac 零模。严格无质量 QED 的带电红外粒子不自动满足孤立极点假设，必须另指定调节或适用的渐近态方案。
电子场与物理单电子态的重叠定义电子极点留数：
\begin{equation}
 \boxed{
 \langle\Omega|\psi(0)|p,s\rangle
 =\sqrt{Z_{\rm pole}^{(e)}}\,u_s(p),}
 \qquad
 (\slashed p-m_{\rm phys}c)u_s(p)=0.
 \label{eq:dirac-field-pole-overlap-dp8}
\end{equation}
于是精确电子二点函数在质量壳附近具有
\begin{equation}
 \boxed{
 \widetilde S(p)
 \xrightarrow[p^2\to m_{\rm phys}^2c^2]{}
 \frac{\ii\hbar Z_{\rm pole}^{(e)}(\slashed p+m_{\rm phys}c)}
 {p^2-m_{\rm phys}^2c^2+\ii0^+}
 +\widetilde S_{\rm reg}(p).}
 \label{eq:dirac-exact-pole-lsz-dp8}
\end{equation}
$\widetilde S_{\rm reg}$ 在该单粒子极点邻域解析，并包含多粒子连续谱贡献。若连通 Green 函数的一条出射电子外腿在该极点附近写成 $\widetilde G(p;\{q\})=\widetilde S(p)\widetilde\Gamma(p;\{q\})+\text{非极点项}$，则电子外腿的留数操作为
\begin{equation}
 \boxed{
 \lim_{p^2\to m_{\rm phys}^2c^2}
 \bar u_s(p)[Z_{\rm pole}^{(e)}]^{-1/2}
 \frac{\slashed p-m_{\rm phys}c}{\ii\hbar}
 \widetilde G(p;\{q\})
 =\sqrt{Z_{\rm pole}^{(e)}}\,\bar u_s(p)\widetilde\Gamma(p;\{q\}).}
 \label{eq:dirac-lsz-amputation-r68}
\end{equation}
\end{proposition}
\begin{proof}
设任意连通 Green 函数含一条出射 Dirac 外腿，写成
\begin{equation*}
 \widetilde G_\alpha(p;\{q\})
 =\int\dd^4 x\,\ee^{+\ii p\cdot x/\hbar}
 \langle\Omega|T\{\psi_\alpha(x)\mathcal O(\{q\})\}|\Omega\rangle_c.
\end{equation*}
若完整谱在 $p^2=m_{\rm phys}^2c^2$ 处含稳定单粒子极点，则极点邻域必可因子化为
\begin{equation*}
 \widetilde G(p;\{q\})
 =\widetilde S(p)\,\widetilde\Gamma(p;\{q\})
 +\widetilde G_{\rm reg}(p;\{q\}),
\end{equation*}
其中 $\widetilde\Gamma$ 与 $\widetilde G_{\rm reg}$ 在该极点解析。正文\eqnrefs{eq:dirac-field-pole-overlap-dp8} 给出场--粒子重叠
\begin{equation*}
 \langle\Omega|\psi(0)|p,s\rangle=\sqrt{Z_{\rm pole}^{(e)}}\,u_s(p),
\end{equation*}
所以精确二点函数的单粒子极点部分为
\begin{equation*}
 \widetilde S(p)
 \xrightarrow[p^2\to m_{\rm phys}^2c^2]{}
 \frac{\ii\hbar Z_{\rm pole}^{(e)}(\slashed p+m_{\rm phys}c)}
 {p^2-m_{\rm phys}^2c^2+\ii0^+}
 +\widetilde S_{\rm reg}(p).
\end{equation*}
左乘自由逆 Dirac 核并使用
\begin{equation*}
 (\slashed p-mc)(\slashed p+mc)=p^2-m^2c^2
\end{equation*}
可知
\begin{equation*}
 [Z_{\rm pole}^{(e)}]^{-1/2}\frac{\slashed p-mc}{\ii\hbar}\,\widetilde S(p)
 \xrightarrow[p^2\to m^2c^2]{}\sqrt{Z_{\rm pole}^{(e)}}
\end{equation*}
在取单粒子极点留数的意义下成立。于是
\begin{equation*}
 \lim_{p^2\to m^2c^2}
 \bar u_s(p)\,
 [Z_{\rm pole}^{(e)}]^{-1/2}\frac{\slashed p-mc}{\ii\hbar}
 \widetilde G(p;\{q\})
 =\sqrt{Z_{\rm pole}^{(e)}}\,\bar u_s(p)\widetilde\Gamma(p;\{q\}),
\end{equation*}
即正文\eqnrefs{eq:dirac-lsz-amputation-r68}。连续谱和其它在该点解析的部分没有相同极点，因此在取留数时消失；严格无质量 QED 的红外粒子结构需按正文给出的调节/包容方案另行处理。
原始截肢核因此还要乘 $\sqrt{Z_{\rm pole}^{(e)}}$，才对应规范化渐近态；其来源是 $Z\cdot Z^{-1/2}=\sqrt Z$。
这里的未投影逆核乘积只是在提取单粒子极点部分的意义下取极限，不能忽略完整矩阵的解析项。
最终与壳上 $\bar u$ 收缩时，逆 Dirac 核的壳上零模消去解析的非极点部分。

\end{proof}


这条式子把“二点函数具有单粒子极点”转换成“外腿可以被截肢并投影到壳上自旋子”。电子与正电子的四种渐近外态只是同一极点投影在交叉变换下的不同取向：
\begin{equation}
 \begin{array}{c|c}
 \text{渐近外态} & \text{壳上波函数}\\
 \hline
 \text{入射 }e^- & u_s(p)\\
 \text{出射 }e^- & \bar u_s(p)\\
 \text{入射 }e^+ & \bar v_s(p)\\
 \text{出射 }e^+ & v_s(p)
 \end{array}
 \label{eq:external-dirac-spinor-table-dp8}
\end{equation}

\begin{proposition}[横向光子单粒子极点的 LSZ 约化]
设物理横向二点函数有单粒子简单极点，$Z_{\rm pole}^{(\gamma)}>0$，原始外腿截肢核在该极点光滑，其余部分不含同阶奇异性。
假定 Ward 恒等式及横向偏振条件成立，先在非零 $k^2$ 计算规范固定逆核，再对物理态取留数。
纵向规范等价类不构成额外渐近粒子。
光子外态满足同样的极点逻辑，但只在两个横向物理偏振上取留数。定义物理光子极点权重
\begin{equation}
 \langle\Omega|A_\mu(0)|k,\lambda\rangle
 =\sqrt{Z_{\rm pole}^{(\gamma)}}\,\varepsilon_\mu^{(\lambda)}(k),
 \qquad
 k^2=0,
 \qquad
 k\cdot\varepsilon^{(\lambda)}=0.
 \label{eq:photon-field-pole-overlap-r68}
\end{equation}
若一条光子外腿的连通 Green 函数写成 $\widetilde G_\nu(k;\{q\})=\widetilde D_{\nu\rho}(k)\widetilde\Gamma^\rho(k;\{q\})+\text{非极点项}$，则相应的留数操作为
\begin{equation}
 \boxed{
 \lim_{k^2\to0}
 \varepsilon_\mu^{(\lambda)*}(k)[Z_{\rm pole}^{(\gamma)}]^{-1/2}
 \frac{\mathscr K_{\xi_{\rm gf}}^{\mu\nu}(k)}{\ii\hbar}
 \widetilde G_\nu(k;\{q\})
 =\sqrt{Z_{\rm pole}^{(\gamma)}}\,\varepsilon_\rho^{(\lambda)*}(k)\widetilde\Gamma^\rho(k;\{q\}).}
 \label{eq:photon-lsz-amputation-r68}
\end{equation}
\end{proposition}
\begin{proof}
物理单光子态 $|k,\lambda\rangle$ 只含两个横向偏振，$\lambda=1,2$。其场--粒子重叠可写成

\begin{equation*}
 \langle\Omega|A_\mu(0)|k,\lambda\rangle
 =\sqrt{Z_{\rm pole}^{(\gamma)}}\,\varepsilon_\mu^{(\lambda)}(k),
 \qquad
 k^2=0,
 \qquad
 k\cdot\varepsilon^{(\lambda)}=0.
\end{equation*}
因此精确光子二点函数在横向单光子极点附近具有
\begin{equation*}
 \widetilde D_{\mu\nu}(k)
 =Z_{\rm pole}^{(\gamma)}\,\widetilde D^{(0)}_{\mu\nu}(k)
 +\widetilde D^{\rm reg}_{\mu\nu}(k),
\end{equation*}
其中 $\widetilde D^{\rm reg}_{\mu\nu}$ 在该极点邻域解析或属于连续谱部分。物理四动量表示的自由逆 Maxwell 核由\eqnrefs{eq:photon-physical-momentum-inverse-kernel-dp42} 定义。结合\eqnrefs{eq:photon-prop-kappa-to-q-r13}，可直接相乘得到
\begin{align*}
 \mathscr K_{\xi_{\rm gf}}^{\mu\rho}(k)
 \widetilde D^{(0)}_{\rho\nu}(k)
 &=\left[-\frac{k^2}{\mu_0\hbar^2}
 \left(P_T^{\mu\rho}+\frac1{\xi_{\rm gf}}P_L^{\mu\rho}\right)\right]\\
 &\quad\times
 \left[-\frac{\ii\mu_0\hbar^3}{k^2+\ii0^+}
 \left(P_{T,\rho\nu}+\xi_{\rm gf}P_{L,\rho\nu}\right)\right],\\
 &\xrightarrow[k^2\to0]{}\ii\hbar\,
 (P_T+P_L)^\mu{}_{\nu},\\
 &=\ii\hbar\,\delta^\mu{}_{\nu}.
\end{align*}
第二行使用 $P_T^2=P_T$、$P_L^2=P_L$ 与 $P_TP_L=0$。于是精确传播子的极点部分满足
\begin{equation*}
 [Z_{\rm pole}^{(\gamma)}]^{-1/2}\frac{\mathscr K_{\xi_{\rm gf}}^{\mu\rho}(k)}{\ii\hbar}
 \widetilde D_{\rho\nu}(k)
 \xrightarrow[k^2\to0]{}
 [Z_{\rm pole}^{(\gamma)}]^{1/2}\delta^\mu{}_{\nu}.
\end{equation*}
再与 $\varepsilon_\mu^{(\lambda)*}$ 或 $\varepsilon_\mu^{(\lambda)}$ 收缩，便得到物理光子外态。因而 $(\ii\hbar)^{-1}$ 不是可省略的约定因子，而是由二点函数的归一化由前述约定固定。

不同协变规范只改变含 $k_\mu$ 或 $k_\nu$ 的纵向部分。物理偏振满足
\begin{equation*}
 k_\mu\varepsilon^{\mu(\lambda)}=0,
\end{equation*}
而与守恒物质流连接时 Ward 恒等式给出
\begin{equation*}
 k_\mu\mathcal M^\mu=0.
\end{equation*}
因此 $\varepsilon^\mu\mapsto\varepsilon^\mu+\alpha k^\mu$ 不改变散射矩阵元；光子外腿表示规范等价类中的横向物理态，而不是四势的四个分量各自对应一个粒子。
故对于 $\widetilde G_\nu=\widetilde D_{\nu\rho}\widetilde\Gamma^\rho+\text{非极点项}$，
物理偏振收缩的留数操作留下 $\sqrt{Z_{\rm pole}^{(\gamma)}}\,\varepsilon_\rho^*\widetilde\Gamma^\rho$。
此平方根需计入 LSZ 归一化核，不能在一般方案中从原始截肢核的定义中省去。

\end{proof}


\begin{proposition}[多外腿的逐极点 LSZ 约化]
设各外部不变量具有以上单粒子极点，抽去总四动量守恒分布后，最高阶多重极点可因子化为完整外腿二点函数乘原始核 $\widetilde\Gamma_{\rm raw}$。
其余项至少缺少一个外腿极点，原始核在相应邻域解析。
令 $\widetilde\Gamma_{\rm amp}$ 为原始核乘所有外腿的 $\sqrt{Z_{\rm pole}}$，逐腿留数及壳上投影给出以下振幅。
规范固定核的纵向部分可以出现在中间传播子中；与物理偏振或守恒流收缩后，它不进入散射矩阵元。对每一条外腿重复同一留数操作，连通 Green 函数的最高阶多重极点便分解为外腿极点与一个共同的截肢核，最终得到
\begin{equation}
 \boxed{
 \mathcal M_{fi}
 =\left[\prod_{\rm ext\;fermions}u,\bar u,v,\bar v\right]
 \left[\prod_{\rm ext\;photons}\varepsilon^{(\lambda)}\right]
 \widetilde\Gamma_{\rm amp}.}
 \label{eq:multi-leg-lsz-r68}
\end{equation}
\end{proposition}
\begin{proof}
考虑含 $N_f$ 条费米子外腿和 $N_\gamma$ 条光子外腿的连通精确 Green 函数。若每一个外部不变量都分别接近一个孤立单粒子极点，多变量 Laurent 展开的最高阶极点部分必可分解为

\begin{equation*}
 \widetilde G_{\rm conn}
 =\left(\prod_{a=1}^{N_f}\widetilde S_a\right)
 \left(\prod_{b=1}^{N_\gamma}\widetilde D_b\right)
 \widetilde\Gamma_{\rm raw}
 +\text{至少缺少一个外腿极点的项}.
\end{equation*}
不同外部坐标具有独立 Fourier 共轭动量，因此逐个外部不变量取留数等价于逐条外腿提取局部 Laurent 系数。对每条费米子外腿乘
\begin{equation*}
 [Z_{\rm pole}^{(e)}]^{-1/2}\frac{\slashed p-mc}{\ii\hbar},
\end{equation*}
对每条光子外腿乘
\begin{equation*}
 [Z_{\rm pole}^{(\gamma)}]^{-1/2}\frac{\mathscr K(k)}{\ii\hbar},
\end{equation*}
再投影到相应壳上自旋子和横向偏振后，外腿传播极点被除去，但每条腿留下极点权重的平方根。定义
\[
 \widetilde\Gamma_{\rm amp}
 =\left(\prod_{a=1}^{N_f}\sqrt{Z_{{\rm pole},a}^{(e)}}\right)
  \left(\prod_{b=1}^{N_\gamma}\sqrt{Z_{{\rm pole},b}^{(\gamma)}}\right)
  \widetilde\Gamma_{\rm raw},
\]
并保持 Grassmann 外腿次序以固定费米交换符号，即得到
\begin{equation*}
 \mathcal M_{fi}
 =\left[\prod_{\text{外部费米子}}\text{壳上 }u,\bar u,v,\bar v\right]
 \left[\prod_{\text{外部光子}}\varepsilon^{(\lambda)}\right]
 \widetilde\Gamma_{\rm amp}.
\end{equation*}
所有相互作用顶点的平移不变性共同产生整体
$(2\pi\hbar)^4\delta^{(4)}(P_f-P_i)$；抽去这一总四动量守恒因子后，$\mathcal M_{fi}$ 才进入截面和衰变率公式。
\end{proof}


因此电子与光子的 LSZ 约化共同实现“精确单粒子极点 $\rightarrow$ 渐近物理态 $\rightarrow$ 连通截肢振幅”的映射。





\begin{derivation}[协变光子模式的对易关系]
在 Feynman 规范 $\xi_{\rm gf}=1$ 中，先把
\[
\mathcal L_{\rm EM,1}
=-\frac1{4\mu_0}F_{\mu\nu}F^{\mu\nu}
-\frac1{2\mu_0}(\partial_\mu A^\mu)^2
\]
展开。利用
\[
F_{\mu\nu}F^{\mu\nu}
=2\bigl[(\partial_\mu A_\nu)(\partial^\mu A^\nu)
-(\partial_\mu A_\nu)(\partial^\nu A^\mu)\bigr]
\]
并对第二项分部积分，交叉项与规范固定项相消；忽略不影响 Euler--Lagrange 方程的全微分后，二次拉氏量变成
\[
\mathcal L_{\rm EM,1}
=-\frac1{2\mu_0}(\partial_\mu A_\nu)(\partial^\mu A^\nu).
\]
因此以 $A^\nu$ 为正则坐标时
\[
\Pi_\nu
\equiv\frac{\partial\mathcal L_{\rm EM,1}}
{\partial(\partial_t A^\nu)}
=-\frac1{\mu_0c^2}\partial_tA_\nu
=-\varepsilon_0\partial_tA_\nu.
\]
这里的下指标十分重要：Minkowski 度规已经进入 $A_\nu=\eta_{\nu\rho}A^\rho$，所以四个分量不会得到同号的正定振子度量。

令
\[
N_{\vb{\kappa}}
\equiv\sqrt{\frac{\hbar}{2\varepsilon_0\omega_{\vb{\kappa}}}}
\]
并把正文\eqnrefs{eq:covariant-photon-fourpol-r7} 写成
\[
A^\mu(x)=\sum_\lambda\int\frac{\dd^3\kappa}{(2\pi)^3}N_{\vb{\kappa}}
\left(\epsilon_\lambda^\mu a_\lambda \ee^{-\ii\kappa\cdot x}
+\epsilon_\lambda^{\mu *}a_\lambda^\dagger \ee^{+\ii\kappa\cdot x}\right).
\]
时间微分给出 $\mp\ii\omega_{\vb{\kappa}}$，所以
\[
\Pi_\nu(x)=\ii\varepsilon_0\sum_\lambda\int\frac{\dd^3\kappa}{(2\pi)^3}
\omega_{\vb{\kappa}}N_{\vb{\kappa}}
\left(\epsilon_{\lambda\nu}a_\lambda \ee^{-\ii\kappa\cdot x}
-\epsilon_{\lambda\nu}^*a_\lambda^\dagger \ee^{+\ii\kappa\cdot x}\right).
\]
暂不预设模代数，而写
\[
[a_\lambda(\vb{\kappa}),a_{\lambda'}^\dagger(\vb{\kappa}')]
=C_{\lambda\lambda'}(\vb{\kappa})
(2\pi)^3\delta^{(3)}(\vb{\kappa}-\vb{\kappa}').
\]
在等时 $x^0=y^0$ 计算 $[A^\mu(\vb{x}),\Pi_\nu(\vb{y})]$。只有 $a$--$a^\dagger$ 与 $a^\dagger$--$a$ 两类收缩；用 $N_{\vb{\kappa}}^2\varepsilon_0\omega_{\vb{\kappa}}=\hbar/2$ 后，两类正负频贡献恰好合成
\begin{align*}
[A^\mu(\vb{x}),\Pi_\nu(\vb{y})]
&=-\ii\hbar\int\frac{\dd^3\kappa}{(2\pi)^3}
 \ee^{\ii\vb{\kappa}\cdot(\vb{x}-\vb{y})}
 \sum_{\lambda\lambda'}C_{\lambda\lambda'}
 \epsilon_\lambda^\mu\epsilon_{\lambda'\nu}^*.
\end{align*}
正则量子化要求这等于
\[
\ii\hbar\delta^\mu{}_{\nu}\delta^{(3)}(\vb{x}-\vb{y})
=\ii\hbar\int\frac{\dd^3\kappa}{(2\pi)^3}
 \ee^{\ii\vb{\kappa}\cdot(\vb{x}-\vb{y})}\delta^\mu{}_{\nu}.
\]
由于四偏振基满足正文\eqnrefs{eq:fourpol-completeness-r7}，即
\[
\sum_{\lambda\lambda'}\eta^{\lambda\lambda'}
\epsilon_\lambda^\mu\epsilon_{\lambda'\nu}^*=\delta^\mu{}_{\nu},
\]
逐个 Fourier 模比较得到唯一解
\[
{C_{\lambda\lambda'}=-\eta_{\lambda\lambda'}}.
\]
于是
\[
[a_\lambda(\vb{\kappa}),a_{\lambda'}^\dagger(\vb{\kappa}')]
=-\eta_{\lambda\lambda'}(2\pi)^3\delta^{(3)}(\vb{\kappa}-\vb{\kappa}'),
\]
即正文\eqnrefs{eq:covariant-photon-mode-comm-r7}。特别地，在这个不定内积空间中，单个时间样激发的自内积为
\[
\langle0|a_0a_0^\dagger|0\rangle<0,
\]
所以负范数并非另加的假设，而是“显式 Lorentz 协变 + 四分量正则量子化 + Minkowski 度规”三者共同固定的代数结果。Gupta--Bleuler 条件随后才负责从这个不定内积空间中选出物理商空间。
\end{derivation}

\begin{derivation}[横向场的正则量子化与模式归一化]
Maxwell 拉氏量

\begin{equation*}
L_{\rm EM}=\frac{\varepsilon_0}{2}\int \dd^3 x\,(\vb{E}^2-c^2\vb{B}^2)
\end{equation*}
给 $\Pi_\phi=0$ 与 $\vb{\Pi}=-\varepsilon_0\vb{E}$。Legendre 变换并对 $\varepsilon_0\vb{E}\cdot\nabla\phi$ 分部积分后
\begin{equation*}
H_{\rm EM}=\int \dd^3 x\left[
\frac{\varepsilon_0}{2}(\vb{E}^2+c^2\vb{B}^2)
+\phi\,\nabla\cdot\vb{\Pi}\right].
\end{equation*}
所以 $\phi$ 是 Lagrange 乘子，变分给 $\nabla\cdot\vb{\Pi}=0$。结合 Coulomb 规范，$A$ 与 $\Pi$ 都只能取横向分量。

写 $P_L=\nabla\nabla/\nabla^2$，则 $P_L^2=P_L$，所以 $P_T=I-P_L$ 满足 $P_T^2=P_T$ 与 $\nabla\cdot P_T=0$。把未施加约束的正则括号 $[A_i,\Pi_j]=\ii\hbar\delta_{ij}\delta^{(3)}$ 两侧各投影一次：
\begin{align*}
[A_i^T(\vb{x}),\Pi_j^T(\vb{y})]
&=(P_T)_{ik}^{(x)}(P_T)_{j\ell}^{(y)}
[A_k(\vb{x}),\Pi_\ell(\vb{y})]\\
&=\ii\hbar(P_T)_{ij}^{(x)}\delta^{(3)}(\vb{x}-\vb{y}),
\end{align*}
这就是\eqnrefs{eq:transverse-comm-r2}。

对固定 $\vb{\kappa}$，$\{\vb{\epsilon}_1,\vb{\epsilon}_2,\hat{\vb{\kappa}}\}$ 构成三维正交归一基，因此单位算符的完备分解给出
\begin{equation*}
\sum_{\lambda=1}^2\epsilon_{\lambda i}\epsilon_{\lambda j}^*
=\delta_{ij}-\hat\kappa_i\hat\kappa_j.
\end{equation*}
再把每个横向 Fourier 系数视为谐振子，以 $[a,a^\dagger]=(2\pi)^3\delta^{(3)}$ 归一化，并要求谐振子哈密顿量为 $\hbar\omega(a^\dagger a+1/2)$，就固定了\eqnrefs{eq:maxwell-mode-r2} 中的 $\sqrt{\hbar/(2\varepsilon_0\omega)}$ 系数。
\end{derivation}

